\subsection*{\texttt{drafts/A045671.txt}}
{\small\raggedright\ttfamily
A045671, section COMMENTS :\par
\par
Kimberling's conjecture -1 \char60{} n*r - a(n) \char60{} 2 is true. More precisely, for n \char62{}= 1, (-9-sqrt(17))/32 \char60{}= n*r - a(n) \char60{}= (27+3*sqrt(17))/32 with r = (1+sqrt(17))/4, and both bounds are best possible (they are the infimum and the supremum). Proof: Dumont-Thomas numeration for the morphism 0-\char62{}11, 1-\char62{}1110 expresses n*r - a(n) as a sum of g\char95{}i*c\char94{}i over the digits, where c = (3-sqrt(17))/2 is the conjugate eigenvalue; the extreme values solve a finite max/min system that is a contraction. - \char95{}Joaquim Manji-Jung\char95{}, Nov xx 2026\par
}
\subsection*{\texttt{drafts/A045672.txt}}
{\small\raggedright\ttfamily
A045672, section COMMENTS :\par
\par
Kimberling's conjecture -1 \char60{} n*r - a(n) \char60{} 3 is true. More precisely, for n \char62{}= 1, (-9-sqrt(17))/16 \char60{}= n*r - a(n) \char60{}= (27+3*sqrt(17))/16 with r = (5+sqrt(17))/2, and both bounds are best possible (they are the infimum and the supremum). Proof: Dumont-Thomas numeration for the morphism 0-\char62{}11, 1-\char62{}1110 expresses n*r - a(n) as a sum of g\char95{}i*c\char94{}i over the digits, where c = (3-sqrt(17))/2 is the conjugate eigenvalue; the extreme values solve a finite max/min system that is a contraction. - \char95{}Joaquim Manji-Jung\char95{}, Nov xx 2026\par
}
\subsection*{\texttt{drafts/A283963.txt}}
{\small\raggedright\ttfamily
A283963, section COMMENTS :\par
\par
Let a(n) be the position of the n-th 0 and b(n) the position of the n-th 1. The conjecture above must be read as -1 \char60{} n*r - a(n) \char60{} 2 and -1 \char60{} n*s - b(n) \char60{} 2 (as literally written, with u(n) the number of 0's \char60{}= n, the quantity n*r - u(n) is unbounded). With this reading the conjecture is true. The exact extreme values are: inf(n*r - a(n)) = (-3-sqrt(3))/6, sup(n*r - a(n)) = (3+sqrt(3))/3, inf(n*s - b(n)) = (-3+sqrt(3))/3, sup(n*s - b(n)) = (3+sqrt(3))/3, with r = (3+sqrt(3))/2, s = sqrt(3). Proof: Dumont-Thomas numeration for the morphism 0-\char62{}1, 1-\char62{}1010 expresses these quantities as sums of g\char95{}i*c\char94{}i over the digits, c being the conjugate eigenvalue; the extreme values solve a finite max/min system that is a contraction. - \char95{}Joaquim Manji-Jung\char95{}, Nov xx 2026\par
}
\subsection*{\texttt{drafts/A283966.txt}}
{\small\raggedright\ttfamily
A283966, section COMMENTS :\par
\par
Let a(n) be the position of the n-th 0 and b(n) the position of the n-th 1. The conjecture above must be read as -1 \char60{} n*r - a(n) \char60{} 3 and -1 \char60{} n*s - b(n) \char60{} 2 (as literally written, with u(n) the number of 0's \char60{}= n, the quantity n*r - u(n) is unbounded). With this reading the conjecture is true. The exact extreme values are: inf(n*r - a(n)) = -1/2, sup(n*r - a(n)) = (5+sqrt(17))/4, inf(n*s - b(n)) = (-7+sqrt(17))/4, sup(n*s - b(n)) = (1+sqrt(17))/4, with r = (7+sqrt(17))/4, s = (-1+sqrt(17))/2. Proof: Dumont-Thomas numeration for the morphism 0-\char62{}1, 1-\char62{}10101 expresses these quantities as sums of g\char95{}i*c\char94{}i over the digits, c being the conjugate eigenvalue; the extreme values solve a finite max/min system that is a contraction. - \char95{}Joaquim Manji-Jung\char95{}, Nov xx 2026\par
}
\subsection*{\texttt{drafts/A284364.txt}}
{\small\raggedright\ttfamily
A284364, section COMMENTS :\par
\par
Let a(n) be the position of the n-th 0 and b(n) the position of the n-th 1. The conjecture above must be read as -2 \char60{} n*r - a(n) \char60{} 2 and -1 \char60{} n*s - b(n) \char60{} 2 (as literally written, with u(n) the number of 0's \char60{}= n, the quantity n*r - u(n) is unbounded). With this reading the conjecture is false. The upper bound 2 fails in both cases: the smallest counterexamples are n = 8933313 (the n-th 0 is at position 20222898) and n = 2977771 (the n-th 1 is at position 5334043), where the value is 2.00487... The exact extreme values are: inf(n*r - a(n)) = (-12-2*sqrt(21))/15, sup(n*r - a(n)) = (6+sqrt(21))/5, inf(n*s - b(n)) = (-9+sqrt(21))/5, sup(n*s - b(n)) = (6+sqrt(21))/5, with r = (9+sqrt(21))/6, s = (-1+sqrt(21))/2. Proof: Dumont-Thomas numeration for the morphism 0-\char62{}1, 1-\char62{}101010 expresses these quantities as sums of g\char95{}i*c\char94{}i over the digits, c being the conjugate eigenvalue; the extreme values solve a finite max/min system that is a contraction. - \char95{}Joaquim Manji-Jung\char95{}, Nov xx 2026\par
}
\subsection*{\texttt{drafts/A284365.txt}}
{\small\raggedright\ttfamily
A284365, section COMMENTS :\par
\par
The conjecture is false: the upper bound 2 fails first at n = 8933313, where a(n) = 20222898 and n*r - a(n) = 2.00487... (exactly: 21*n\char94{}2 \char62{} (6*a(n)+12-9*n)\char94{}2). The exact extreme values are inf(n*r - a(n)) = (-12-2*sqrt(21))/15 = -1.411... and sup(n*r - a(n)) = (6+sqrt(21))/5 = 2.1165..., neither being attained. See A284364 and the link for the method (Dumont-Thomas numeration for the morphism 0-\char62{}1, 1-\char62{}101010). - \char95{}Joaquim Manji-Jung\char95{}, Nov xx 2026\par
}
\subsection*{\texttt{drafts/A284366.txt}}
{\small\raggedright\ttfamily
A284366, section COMMENTS :\par
\par
The conjecture is false: the upper bound 2 fails first at n = 2977771, where a(n) = 5334043 and n*s - a(n) = 2.00487... (exactly: 21*n\char94{}2 \char62{} (2*a(n)+4+n)\char94{}2). The lower bound -1 holds. The exact extreme values are inf(n*s - a(n)) = (-9+sqrt(21))/5 = -0.8834... and sup(n*s - a(n)) = (6+sqrt(21))/5 = 2.1165..., neither being attained. See A284364 and the link for the method (Dumont-Thomas numeration for the morphism 0-\char62{}1, 1-\char62{}101010). - \char95{}Joaquim Manji-Jung\char95{}, Nov xx 2026\par
}
\subsection*{\texttt{drafts/A284368.txt}}
{\small\raggedright\ttfamily
A284368, section COMMENTS :\par
\par
Let a(n) be the position of the n-th 0 and b(n) the position of the n-th 1. The conjecture above must be read as 1 \char60{} n*r - a(n) \char60{} 3 and -1 \char60{} n*s - b(n) \char60{} 1 (as literally written, with u(n) the number of 0's \char60{}= n, the quantity n*r - u(n) is unbounded). With this reading the conjecture is true. The exact extreme values are: inf(n*r - a(n)) = (1+sqrt(13))/3, sup(n*r - a(n)) = (4+sqrt(13))/3, inf(n*s - b(n)) = (-5+sqrt(13))/3, sup(n*s - b(n)) = (-2+sqrt(13))/3, with r = (5+sqrt(13))/2, s = (-1+sqrt(13))/2. Proof: Dumont-Thomas numeration for the morphism 0-\char62{}1, 1-\char62{}1011 expresses these quantities as sums of g\char95{}i*c\char94{}i over the digits, c being the conjugate eigenvalue; the extreme values solve a finite max/min system that is a contraction. - \char95{}Joaquim Manji-Jung\char95{}, Nov xx 2026\par
}
\subsection*{\texttt{drafts/A284369.txt}}
{\small\raggedright\ttfamily
A284369, section COMMENTS :\par
\par
Let a(n) be the position of the n-th 0 and b(n) the position of the n-th 1. The conjecture above must be read as -3 \char60{} n*r - a(n) \char60{} 4 and -2 \char60{} n*s - b(n) \char60{} 3 (as literally written, with u(n) the number of 0's \char60{}= n, the quantity n*r - u(n) is unbounded). With this reading the conjecture is true. The bound -2 for the 1's is the exact infimum; it is not attained since n*sqrt(3) is irrational. The exact extreme values are: inf(n*r - a(n)) = (-3-sqrt(3))/2, sup(n*r - a(n)) = 2+sqrt(3), inf(n*s - b(n)) = -2, sup(n*s - b(n)) = 1+sqrt(3), with r = (3+sqrt(3))/2, s = sqrt(3). Proof: Dumont-Thomas numeration for the morphism 0-\char62{}1, 1-\char62{}1001 expresses these quantities as sums of g\char95{}i*c\char94{}i over the digits, c being the conjugate eigenvalue; the extreme values solve a finite max/min system that is a contraction. - \char95{}Joaquim Manji-Jung\char95{}, Nov xx 2026\par
}
\subsection*{\texttt{drafts/A284386.txt}}
{\small\raggedright\ttfamily
A284386, section COMMENTS :\par
\par
Let a(n) be the position of the n-th 0 and b(n) the position of the n-th 1. The conjecture above must be read as 0 \char60{} n*r - a(n) \char60{} 2 and -1 \char60{} n*s - b(n) \char60{} 1 (as literally written, with u(n) the number of 0's \char60{}= n, the quantity n*r - u(n) is unbounded). With this reading the conjecture is true. Here r must be read as (5+sqrt(13))/2; the value (5+sqrt(3))/2 in the comment is a typo (it is not 1/(1-1/s)). The exact extreme values are: inf(n*r - a(n)) = (1+sqrt(13))/6, sup(n*r - a(n)) = (7+sqrt(13))/6, inf(n*s - b(n)) = (-5+sqrt(13))/6, sup(n*s - b(n)) = (1+sqrt(13))/6, with r = (5+sqrt(13))/2, s = (-1+sqrt(13))/2. Proof: Dumont-Thomas numeration for the morphism 0-\char62{}1, 1-\char62{}1101 expresses these quantities as sums of g\char95{}i*c\char94{}i over the digits, c being the conjugate eigenvalue; the extreme values solve a finite max/min system that is a contraction. - \char95{}Joaquim Manji-Jung\char95{}, Nov xx 2026\par
}
\subsection*{\texttt{drafts/A284505.txt}}
{\small\raggedright\ttfamily
A284505, section COMMENTS :\par
\par
Let a(n) be the position of the n-th 0 and b(n) the position of the n-th 1. The conjecture above must be read as -3 \char60{} n*r - a(n) \char60{} 4 and -2 \char60{} n*s - b(n) \char60{} 4 (as literally written, with u(n) the number of 0's \char60{}= n, the quantity n*r - u(n) is unbounded). With this reading the conjecture is true. The exact extreme values are: inf(n*r - a(n)) = (-9-5*sqrt(3))/6, sup(n*r - a(n)) = (6+2*sqrt(3))/3, inf(n*s - b(n)) = (-3-sqrt(3))/3, sup(n*s - b(n)) = (6+2*sqrt(3))/3, with r = (3+sqrt(3))/2, s = sqrt(3). Proof: Dumont-Thomas numeration for the morphism 0-\char62{}1, 1-\char62{}1100 expresses these quantities as sums of g\char95{}i*c\char94{}i over the digits, c being the conjugate eigenvalue; the extreme values solve a finite max/min system that is a contraction. - \char95{}Joaquim Manji-Jung\char95{}, Nov xx 2026\par
}
\subsection*{\texttt{drafts/A284677.txt}}
{\small\raggedright\ttfamily
A284677, section FORMULA :\par
\par
The conjecture a(n) = A284674(n-1) fails at n = 6: a(6) = 0 and A284674(5) = 1. The correct relation is a(n) = A284674(n) for n \char62{}= 3. Proof: let p: 0-\char62{}1, 1-\char62{}0111; then p maps each limiting word to the other. Writing a = 10t and A284674 = 01u gives u = 111p(t) and t = 111p(u), and the map w -\char62{} 111p(111p(w)) has a unique infinite fixed point, so t = u. - \char95{}Joaquim Manji-Jung\char95{}, Nov xx 2026\par
}
\subsection*{\texttt{drafts/A284948.txt}}
{\small\raggedright\ttfamily
A284948, section FORMULA (ou COMMENTS) :\par
\par
The conjecture a(n) = A096268(n-2), n\char62{}1, is true. Proof: both sequences are 4-automatic (this one is the fixed point starting with 1 of 0-\char62{}0010, 1-\char62{}1010, and A096268 is the fixed point of 0-\char62{}0100, 1-\char62{}0101), and the statement "for all i\char62{}=0, a(i+2) = A096268(i)" is verified by the Walnut theorem prover. - \char95{}Joaquim Manji-Jung\char95{}, Nov xx 2026\par
}
\subsection*{\texttt{drafts/A285076.txt}}
{\small\raggedright\ttfamily
A285076, section FORMULA :\par
\par
The conjecture a(n) = A285073(n), n\char62{}=3, is true; it follows from the comments by Michel Dekking above (a = 10c and A285073 = 01c with the same word c = A159684). - \char95{}Joaquim Manji-Jung\char95{}, Nov xx 2026\par
}
\subsection*{\texttt{drafts/A285358.txt}}
{\small\raggedright\ttfamily
A285358, section FORMULA :\par
\par
The conjecture a(n) = A284929(n+1) is true. Proof: let s: 0-\char62{}10, 1-\char62{}1101 and t: 0-\char62{}01, 1-\char62{}1011. Then s(w)1 = 1t(w) for every finite word w, so the fixed point y = A284929 of t satisfies s(y) = 1y. Writing y = 0z gives 10s(z) = 10z, so z is the fixed point of s starting with 1. - \char95{}Joaquim Manji-Jung\char95{}, Nov xx 2026\par
}
\subsection*{\texttt{conjectures/A000139/draft.txt}}
{\small\raggedright\ttfamily
\%C A000139 Bala's conjecture is true. By Legendre's formula, v\char95{}2(a(n)) = s(n) + s(n+1) - s(3n) = c + 1 - t, where s = A000120, t is the number of trailing 1's of n and c is the number of carries in the binary addition n + 2n. If t \char62{}= 2 then c \char62{}= t; so a(n) is odd iff t = 1 and c = 0, i.e., iff n == 1 (mod 4) and n has no two adjacent 1's in binary, i.e., n is in A022341. - \char95{}Joaquim Manji-Jung\char95{}, Nov 18 2026\par
}
\subsection*{\texttt{conjectures/A001462/draft.txt}}
{\small\raggedright\ttfamily
\%C A001462 The conjecture a(n) \char62{}= n\char94{}(phi-1) of Jianing Song is true. Let S(m) = A001463(m), the last position of m; since phi*(phi-1) = 1 it is equivalent to S(m) \char60{}= m\char94{}phi. One proves m\char94{}phi/phi \char60{}= S(m) \char60{}= m\char94{}phi by strong induction: if a(j) = v then S(v-1) \char60{} j \char60{}= S(v), which gives j\char94{}(phi-1) \char60{}= v \char60{} 1 + (phi*j)\char94{}(phi-1); summing over j \char60{}= m gives the two bounds for m \char62{} 30, and m \char60{}= 30 is checked directly. - \char95{}Joaquim Manji-Jung\char95{}, Nov 18 2026\par
}
\subsection*{\texttt{conjectures/A003278/draft.txt}}
{\small\raggedright\ttfamily
\%C A003278 Jeffery's conjecture is true: a(n) - 1 = A005836(n), A191107(n) = 1 + 3*A005836(n) and A055246(n) = 1 + 6*A005836(n) (see A055246). - \char95{}Joaquim Manji-Jung\char95{}, Nov 18 2026\par
}
\subsection*{\texttt{conjectures/A003959/draft.txt}}
{\small\raggedright\ttfamily
\%C A003959 Stephan's formula is true: a is multiplicative with a(p\char94{}e) = (p+1)\char94{}e, so the Euler factor at p is Sum\char95{}\char123{}e\char62{}=0\char125{} ((p+1)/p\char94{}s)\char94{}e = 1/(1 - p\char94{}(-s) - p\char94{}(1-s)) (for Re(s) \char62{} 2). - \char95{}Joaquim Manji-Jung\char95{}, Nov 19 2026\par
}
\subsection*{\texttt{conjectures/A004123/draft.txt}}
{\small\raggedright\ttfamily
\%C A004123 Bala's conjectures are true, in the following form: if a(n) has e.g.f. G(exp(x)-1) with G an integral power series, then a(n) = Sum\char95{}j g\char95{}j*j!*Stirling2(n,j), and since k divides j! for j \char62{}= k, a(n) == Sum\char95{}\char123{}i\char60{}k\char125{} c\char95{}i*i\char94{}n (mod k) for some integers c\char95{}i. Hence a(n) mod k is periodic for n \char62{}= log\char95{}2(k), with period dividing the Carmichael function A002322(k), which divides phi(k). - \char95{}Joaquim Manji-Jung\char95{}, Nov 18 2026\par
\%C A006531 The conjecture is true: Sum\char95{}k (-1)\char94{}k*Catalan(k)*k!*Stirling2(n,k) == Sum\char95{}\char123{}i\char60{}m\char125{} c\char95{}i*i\char94{}n (mod m), since m divides k! for k \char62{}= m, so a(n) mod m is periodic for n \char62{}= log\char95{}2(m) with period dividing A002322(m), which divides phi(m) (the factor (-1)\char94{}n has period 2). See A004123. - \char95{}Joaquim Manji-Jung\char95{}, Nov 18 2026\par
}
\subsection*{\texttt{conjectures/A005121/draft.txt}}
{\small\raggedright\ttfamily
\%C A005121 Mohapatra's conjecture is true. By induction R(n,k) = Sum\char95{}\char123{}m\char60{}n\char125{} c(n,k,m)*R(m,m), where for fixed m the vector (c(n,k,m))\char95{}k equals M\char94{}(n-m-1) applied to (1,2,...,m), with M[k][j] = k*[j \char62{}= k]; the row vector 1*M\char94{}t has entries Stirling2(j+t,j) (same recurrence), so Sum\char95{}k c(n,k,m) = Stirling2(n,m) by telescoping, and R(n,n) satisfies Lengyel's recurrence a(n) = Sum\char95{}\char123{}m\char60{}n\char125{} Stirling2(n,m)*a(m). - \char95{}Joaquim Manji-Jung\char95{}, Nov 19 2026\par
}
\subsection*{\texttt{conjectures/A005427/draft.txt}}
{\small\raggedright\ttfamily
\%C A005427 Hadjicostas's general conjecture (Jul 22 2020) is true for every k \char62{}= 2. Writing r\char95{}n = (k-1)*T(n;k) - Sum\char95{}\char123{}s\char60{}n\char125{} T(s;k) (in [0, k-2]), one has J\char95{}k(T(n;k)) = r\char95{}n + 1 for T(n;k) \char62{}= k, where J\char95{}k(m) is the survivor: for m \char62{}= k-1 at most one wrap-around happens per step, so starting from J(m0) = j \char60{}= k-1 the survivor is j + k*s \char62{}= k+1 until the first s with j + (k-1)*s \char62{} m0, which is s = ceiling((m0 - r\char95{}n)/(k-1)) = T(n+1;k) - T(n;k), where it becomes j + (k-1)*s - m0 = r\char95{}\char123{}n+1\char125{} + 1. The base case m = k follows from J\char95{}k(v) = k*v + 1 - Q(v) for v \char60{}= k-1, Q(v) being the partial sum of T at the last occurrence of v. - \char95{}Joaquim Manji-Jung\char95{}, Nov 19 2026\par
}
\subsection*{\texttt{conjectures/A005590/draft.txt}}
{\small\raggedright\ttfamily
\%C A005590 Stephan's conjecture is true: a(3n) = 0 if and only if n is in A003714, and otherwise a(3n) is even with |a(3n)| \char62{}= 2. Proof: the vector u(n) = (a(3n), a(3n+1), a(3n+2), a(3n+3)) satisfies u(2n) = (a0, a1-a0, a1, a2-a1) and u(2n+1) = (a2-a1, a2, a3-a2, a3) in terms of u(n) = (a0, a1, a2, a3). While the binary digits of n contain no 11, u(n) is (0, 1, 1, -2m) or (0, 1, -2m-1, -2m); after the first 11 it lies in the set \char123{}a0(a0-a1) \char62{} 0, a1(a1-a2) \char62{} 0, a2(a2-a3) \char62{} 0\char125{}, which is mapped into itself by both maps and excludes a0 = 0. - \char95{}Joaquim Manji-Jung\char95{}, Nov 18 2026\par
}
\subsection*{\texttt{conjectures/A006068/draft.txt}}
{\small\raggedright\ttfamily
\%C A006068 Kurkov's conjecture is true: the inverse Gray code is a synchronized function in base 2 (bit i of a(n) is the XOR of the bits of n in positions \char62{}= i), and the identity is proved with the Walnut theorem prover. - \char95{}Joaquim Manji-Jung\char95{}, Nov 19 2026\par
}
\subsection*{\texttt{conjectures/A006165/draft.txt}}
{\small\raggedright\ttfamily
\%C A006165 The conjectures (i)-(vii) of Peter Bala are true. With p the leading power of 2 of n \char62{}= 2, a(n) = min(n - p/2, p), and by induction a\char94{}(k)(n) = min(n - (2\char94{}k-1)*p/2\char94{}k, 2p/2\char94{}k) for n \char62{}= 2\char94{}k, a\char94{}(k)(n) = 1 for n \char60{} 2\char94{}k. Then (i)-(v) are proved with the Walnut theorem prover, and (vi) for all k \char62{}= 2 at once, by taking u = p/2\char94{}k as a second variable (it fails for k = 1, e.g., at n = 5). For (vii), n - a\char94{}(k)(n) = max((2\char94{}k-1)*p/2\char94{}k, n - 2p/2\char94{}k), so the repeated values are (2\char94{}k-1)*2\char94{}j, each repeated 2\char94{}j + 1 times (A000051: 2, 3, 5, 9, 17, 33, ...; the 35 above should be 33). - \char95{}Joaquim Manji-Jung\char95{}, Nov 18 2026\par
}
\subsection*{\texttt{conjectures/A006364/draft.txt}}
{\small\raggedright\ttfamily
\%C A006364 Yanev's formula is true: the terms come in pairs 2j, 2j+1 with j evil, and the i-th evil number (from i = 0) is 2i + A010060(i), so a(2i+1) = 4i + 2*A010060(i) and a(2i+2) = a(2i+1) + 1, which is the formula. - \char95{}Joaquim Manji-Jung\char95{}, Nov 18 2026\par
}
\subsection*{\texttt{conjectures/A007226/draft.txt}}
{\small\raggedright\ttfamily
\%C A007226 Bala's conjecture is true, since a(n) = (2n+1)*A000139(n) and A000139(n) is odd iff n is in A022341 (see there). - \char95{}Joaquim Manji-Jung\char95{}, Nov 19 2026\par
}
\subsection*{\texttt{conjectures/A007413/draft.txt}}
{\small\raggedright\ttfamily
\%C A007413 Shevelev's conjecture is true. Since A076826(2m+1) = 1 and A076826(2m) = 2*A010060(m), it is equivalent to a(n) = A000069(n+1) - A000069(n), and with A000069(j+1) = 2j + 1 - A010060(j) one checks with Walnut that these differences are the fixed point of 1 -\char62{} 123, 2 -\char62{} 13, 3 -\char62{} 2 (the block of the j-th letter starts at position 2j + A010060(j)). - \char95{}Joaquim Manji-Jung\char95{}, Nov 18 2026\par
}
\subsection*{\texttt{conjectures/A008278/draft.txt}}
{\small\raggedright\ttfamily
\%C A008278 Wiseman's conjecture is true: such a divisor amounts to placing k-1 non-attacking rooks on the staircase board with rows of lengths 1, 2, ..., n-1 (prime p\char95{}i has exponent at most n-i), and the rook numbers of this board satisfy r(n,j) = r(n-1,j) + (n-j)*r(n-1,j-1), the recurrence of Stirling2(n,n-j). - \char95{}Joaquim Manji-Jung\char95{}, Nov 19 2026\par
}
\subsection*{\texttt{conjectures/A014664/draft.txt}}
{\small\raggedright\ttfamily
\%C A014664 Shevelev's conjecture is true (classical): if d = a(n) and 2\char94{}d = 1 + p*t with p not dividing t (equivalent to p not being a Wieferich prime), the lifting-the-exponent lemma gives v\char95{}p(2\char94{}(d*p\char94{}i) - 1) = 1 + i, so the order of 2 modulo p\char94{}k is d*p\char94{}(k-1). - \char95{}Joaquim Manji-Jung\char95{}, Nov 19 2026\par
}
\subsection*{\texttt{conjectures/A022344/draft.txt}}
{\small\raggedright\ttfamily
\%C A022344 Ramsey's conjecture is true: F3 is a Fibonacci sequence (linear in F1), and with x = F3(i-1), y = F3(i) its characteristic is |y\char94{}2 - x*y - x\char94{}2|, the norm form of Z[phi], which gives |Q(F3)| = |Q(F1)|*|Q(F2)| (an identity in the initial terms); F3 is nonzero and of constant sign from some point on, so by Morrison's theorem |F3| eventually coincides with a row p of the Wythoff array, and a(p) = a(n)*a(m). - \char95{}Joaquim Manji-Jung\char95{}, Nov 19 2026\par
}
\subsection*{\texttt{conjectures/A033485/draft.txt}}
{\small\raggedright\ttfamily
\%C A033485 The conjectured value c = 1.64... is wrong: the limit exists and equals 2*log(2) = 1.3862... Writing beta(n) = n*a(floor(n/2))/a(n), one has a(2N) = 2*S(N) - a(N) + 1 with S(N) = a(1)+...+a(N), and S(N)/a(N) = Sum\char95{}k Prod\char95{}\char123{}j\char62{}k\char125{} (1 - beta(j)/j); if |beta(j) - log\char95{}2(j)| \char60{}= D on (N/2, N] this gives |beta(n) - log\char95{}2(n)| \char60{}= D + O(log log N/log N) for n = 2N, 2N+1, and a dyadic induction (base computed up to 2\char94{}29) yields beta(n) = log\char95{}2(n) + O((log log n)\char94{}2), hence a(2n)/a(n)*log(n)/n = 2*log(n)/beta(2n) -\char62{} 2*log(2). The ratio is 1.566 at n = 1.5*10\char94{}8 and still decreasing. - \char95{}Joaquim Manji-Jung\char95{}, Nov 19 2026\par
}
\subsection*{\texttt{conjectures/A034448/draft.txt}}
{\small\raggedright\ttfamily
\%C A034448 Yanev's conjecture is true: both sides are multiplicative, and for n = p\char94{}e, sigma(p\char94{}(2e-1))/sigma(p\char94{}(e-1)) = (p\char94{}(2e)-1)/(p\char94{}e-1) = p\char94{}e + 1. - \char95{}Joaquim Manji-Jung\char95{}, Nov 18 2026\par
}
\subsection*{\texttt{conjectures/A036766/draft.txt}}
{\small\raggedright\ttfamily
\%C A036766 The conjectured recurrence is true (for n \char62{}= 2): with theta = x d/dx, the recurrence corresponds to the operator L = Sum\char95{}i x\char94{}i p\char95{}i(theta+i), and reducing L(A) modulo A = Sum\char95{}\char123{}k=0..4\char125{} (x*A)\char94{}k (using A' = -P\char95{}x/P\char95{}A) gives L(A) = 8328 + 6660*x. - \char95{}Joaquim Manji-Jung\char95{}, Nov 19 2026\par
}
\subsection*{\texttt{conjectures/A040027/draft.txt}}
{\small\raggedright\ttfamily
\%C A040027 Yanev's conjecture is true: Sum\char95{}n Stirling2(n+1,k+1)*x\char94{}n/n! = exp(x)*(exp(x)-1)\char94{}k/k!, so the right side has e.g.f. H(x) = exp(x)*C(exp(x)-1) with C(y) = Sum\char95{}k A058006(k)*y\char94{}k/k!; since C' - C = -1/(1+y)\char94{}2, H' = (1+exp(x))*H - 1 with H(0) = 1, which is the derivative form of the e.g.f. equation above (for the sequence shifted by one). - \char95{}Joaquim Manji-Jung\char95{}, Nov 19 2026\par
}
\subsection*{\texttt{conjectures/A050446/draft.txt}}
{\small\raggedright\ttfamily
\%C A050446 Rascoe's conjecture is true: replacing each letter x in even position by m+1-x maps these words bijectively onto the words with w1 \char60{}= w2 \char62{}= w3 \char60{}= ... over [m], and splitting a word w1 \char62{}= w2 \char60{}= w3 \char62{}= ... at its first letter m (which is in odd position 2k+1) gives the recurrence T(n,m) = T(n,m-1) + Sum\char95{}k T(2k,m-1)*T(n-1-2k,m) above. - \char95{}Joaquim Manji-Jung\char95{}, Nov 18 2026\par
}
\subsection*{\texttt{conjectures/A055246/draft.txt}}
{\small\raggedright\ttfamily
\%C A055246 Both conjectures of Jeffery are true and follow from a(n) = 1 + 6*A005836(n): A005836 is generated by 0, c -\char62{} 3c, c -\char62{} 3c+1, and 3*(1+6c) - 2 = 1 + 6*(3c), 3*(1+6c) + 4 = 1 + 6*(3c+1); A191107(n) = 1 + 3*A005836(n) and A003278(n) = A005836(n) + 1; and the binary expansion of (2n-1)*2\char94{}(k-1) - 1 is that of n-1 followed by 0 and k-1 ones, so A005836((2n-1)*2\char94{}(k-1)) = 3\char94{}k*A005836(n) + (3\char94{}(k-1)-1)/2. - \char95{}Joaquim Manji-Jung\char95{}, Nov 18 2026\par
}
\subsection*{\texttt{conjectures/A055938/draft.txt}}
{\small\raggedright\ttfamily
\%C A055938 Stephan's conjectures are true. Since A005187(m+1) - A005187(m) = 1 + A007814(m+1), the terms come in runs of consecutive integers separated by gaps of 3, and the gaps of 3 occur at the indices n = A005187(j), j \char62{}= 1; hence a(n+1) - a(n) = 1 iff n is in the sequence. The identity a(n) = A080578(n+1) - 2 then follows by induction, since A080578 obeys the same rule. - \char95{}Joaquim Manji-Jung\char95{}, Nov 19 2026\par
}
\subsection*{\texttt{conjectures/A060973/draft.txt}}
{\small\raggedright\ttfamily
\%C A060973 Both conjectures of Peter Bala are true; since a(n) = n - A006165(n) has a closed form synchronized in base 2, they are proved with the Walnut theorem prover. - \char95{}Joaquim Manji-Jung\char95{}, Nov 19 2026\par
}
\subsection*{\texttt{conjectures/A070864/draft.txt}}
{\small\raggedright\ttfamily
\%C A070864 Layman's conjecture is true; it is equivalent to the formulas of Cloitre and Curgus: a(k\char94{}2) = 2k-1 and a(k\char94{}2-i) = 2k-2-(-1)\char94{}i for 1 \char60{}= i \char60{}= 2k-2. Proof by strong induction: for n = k\char94{}2-i with 1 \char60{}= i \char60{}= 2k-3 one has a(n-1) = 2k-2+(-1)\char94{}i, so n - a(n-1) = (k-1)\char94{}2 - i' with i' = i-1+(-1)\char94{}i, which lies in the previous block; the cases n = k\char94{}2 and n = (k-1)\char94{}2+1 are checked in the same way. - \char95{}Joaquim Manji-Jung\char95{}, Nov 18 2026\par
}
\subsection*{\texttt{conjectures/A072939/draft.txt}}
{\small\raggedright\ttfamily
\%C A072939 All four conjectures are true. Writing d = c(k+1) - c(k): an even d becomes -d/2; an odd negative d stays and the terms increase by |d| every two steps; an odd positive d stays and the terms decrease by d while d \char60{} c(k). For odd n with n - 1 = 2\char94{}v*u, u odd, the first v steps give d = (-1)\char94{}v*u, so c(k) -\char62{} oo iff v is odd; if v is even the pair reached is (A, A+u) with A == 1 (mod u), which leads to (1, 1+u) and then to the constant 1 (as for even n). Hence a(n) = A036554(n) + 1 and a(n)/n -\char62{} 3. Since A328979 satisfies x(4i+3) = 1, x(4i+2) = x(4i+4) = 0, x(4i+1) = x(i+1) (offset 1), its 1's are at the positions A036554(n) + 1, and A285384 is its complement. - \char95{}Joaquim Manji-Jung\char95{}, Nov 19 2026\par
}
\subsection*{\texttt{conjectures/A074355/draft.txt}}
{\small\raggedright\ttfamily
\%C A074355 The conjectured g.f. is false as written (its x\char94{}4 coefficient is 27, not 15); the correct g.f. is (9x\char94{}4+3x\char94{}3)/(1-x-3x\char94{}2)\char94{}2, which follows from [q\char94{}1]: a(n) = a(n-1) + 3a(n-2) + 3*A006130(n-2) for n \char62{}= 3. - \char95{}Joaquim Manji-Jung\char95{}, Nov 19 2026\par
}
\subsection*{\texttt{conjectures/A074357/draft.txt}}
{\small\raggedright\ttfamily
\%C A074356 The conjectured g.f. is true: [q\char94{}i](1+...+q\char94{}(n-2)) = [i \char60{}= n-2] gives c\char95{}j(n) = b*c\char95{}j(n-1) + lambda*Sum\char95{}\char123{}i=0..min(j,n-2)\char125{} c\char95{}\char123{}j-i\char125{}(n-2), hence C\char95{}j(x)*(1 - b*x - lambda*x\char94{}2) = lambda*x\char94{}2*Sum\char95{}\char123{}i=1..j\char125{} (C\char95{}\char123{}j-i\char125{}(x) - (terms of degree \char60{} i)), a triangular system whose solution is the stated rational function. - \char95{}Joaquim Manji-Jung\char95{}, Nov 19 2026\par
\%C A074357 The conjectured g.f. is true: [q\char94{}i](1+...+q\char94{}(n-2)) = [i \char60{}= n-2] gives c\char95{}j(n) = b*c\char95{}j(n-1) + lambda*Sum\char95{}\char123{}i=0..min(j,n-2)\char125{} c\char95{}\char123{}j-i\char125{}(n-2), hence C\char95{}j(x)*(1 - b*x - lambda*x\char94{}2) = lambda*x\char94{}2*Sum\char95{}\char123{}i=1..j\char125{} (C\char95{}\char123{}j-i\char125{}(x) - (terms of degree \char60{} i)), a triangular system whose solution is the stated rational function. - \char95{}Joaquim Manji-Jung\char95{}, Nov 19 2026\par
\%C A074358 The conjectured g.f. is true: [q\char94{}i](1+...+q\char94{}(n-2)) = [i \char60{}= n-2] gives c\char95{}j(n) = b*c\char95{}j(n-1) + lambda*Sum\char95{}\char123{}i=0..min(j,n-2)\char125{} c\char95{}\char123{}j-i\char125{}(n-2), hence C\char95{}j(x)*(1 - b*x - lambda*x\char94{}2) = lambda*x\char94{}2*Sum\char95{}\char123{}i=1..j\char125{} (C\char95{}\char123{}j-i\char125{}(x) - (terms of degree \char60{} i)), a triangular system whose solution is the stated rational function. - \char95{}Joaquim Manji-Jung\char95{}, Nov 19 2026\par
\%C A074359 The conjectured g.f. is true: [q\char94{}i](1+...+q\char94{}(n-2)) = [i \char60{}= n-2] gives c\char95{}j(n) = b*c\char95{}j(n-1) + lambda*Sum\char95{}\char123{}i=0..min(j,n-2)\char125{} c\char95{}\char123{}j-i\char125{}(n-2), hence C\char95{}j(x)*(1 - b*x - lambda*x\char94{}2) = lambda*x\char94{}2*Sum\char95{}\char123{}i=1..j\char125{} (C\char95{}\char123{}j-i\char125{}(x) - (terms of degree \char60{} i)), a triangular system whose solution is the stated rational function. - \char95{}Joaquim Manji-Jung\char95{}, Nov 19 2026\par
\%C A074360 The conjectured g.f. is true: [q\char94{}i](1+...+q\char94{}(n-2)) = [i \char60{}= n-2] gives c\char95{}j(n) = b*c\char95{}j(n-1) + lambda*Sum\char95{}\char123{}i=0..min(j,n-2)\char125{} c\char95{}\char123{}j-i\char125{}(n-2), hence C\char95{}j(x)*(1 - b*x - lambda*x\char94{}2) = lambda*x\char94{}2*Sum\char95{}\char123{}i=1..j\char125{} (C\char95{}\char123{}j-i\char125{}(x) - (terms of degree \char60{} i)), a triangular system whose solution is the stated rational function. - \char95{}Joaquim Manji-Jung\char95{}, Nov 19 2026\par
}
\subsection*{\texttt{conjectures/A078841/draft.txt}}
{\small\raggedright\ttfamily
\%C A078841 The conjecture is true: writing m = 2\char94{}(n-j)*u with u odd and Omega(u) = j, the map m -\char62{} m/2\char94{}n is an increasing bijection from the n-almost primes onto \char123{}u/2\char94{}Omega(u): u odd, Omega(u) \char60{}= n\char125{}, whose n+1 smallest elements do not depend on n (elements with Omega(u) = j are \char62{}= (3/2)\char94{}j). Hence a(n) = 2\char94{}n*t(n+1), where t(1) \char60{}= t(2) \char60{}= ... enumerates V = \char123{}u/2\char94{}Omega(u): u odd\char125{}; since V contains p/2 for every odd prime p, the prime number theorem gives t(k+1)/t(k) -\char62{} 1, so a(n+1)/a(n) -\char62{} 2. - \char95{}Joaquim Manji-Jung\char95{}, Nov 19 2026\par
}
\subsection*{\texttt{conjectures/A079267/draft.txt}}
{\small\raggedright\ttfamily
\%C A079267 Callan's conjecture is true: by the explicit formula, d(n,s)/(2n-1)!! = (1/s!)*Sum\char95{}\char123{}h\char62{}=s\char125{} (-1)\char94{}(h-s)*r\char95{}h/(h-s)! with r\char95{}h = Product\char95{}\char123{}i\char60{}h\char125{} 2(n-i)/(2n-i), so 0 \char60{}= r\char95{}h \char60{}= 1 and r\char95{}h -\char62{} 1 for fixed h; by dominated convergence the limit is exp(-1)/s!. - \char95{}Joaquim Manji-Jung\char95{}, Nov 19 2026\par
}
\subsection*{\texttt{conjectures/A080253/draft.txt}}
{\small\raggedright\ttfamily
\%C A080253 Bala's conjecture is true: exp(x)/(2-exp(2x)) = G(exp(x)-1) with G(u) = (1+u)/(1-2u-u\char94{}2) an integral power series, so a(n) = Sum\char95{}j g\char95{}j*j!*Stirling2(n,j) and, as in A004123, a(n) mod k is periodic for n \char62{}= log\char95{}2(k) with period dividing A002322(k), hence phi(k). - \char95{}Joaquim Manji-Jung\char95{}, Nov 19 2026\par
}
\subsection*{\texttt{conjectures/A083207/draft.txt}}
{\small\raggedright\ttfamily
\%C A083207 The conjecture that any 12 consecutive numbers include a Zumkeller number is true: 3*2\char94{}a is Zumkeller for a \char62{}= 1 (take \char123{}3*2\char94{}a, 2, 4, ..., 2\char94{}(a-1)\char125{}), and q*m is Zumkeller whenever q is and gcd(q,m) = 1 (multiply both halves by the divisors of m); among 12 consecutive numbers there are two consecutive multiples of 6, one of which has 3-adic valuation 1, hence is 3*2\char94{}a*m with gcd(m,6) = 1. - \char95{}Joaquim Manji-Jung\char95{}, Nov 19 2026\par
}
\subsection*{\texttt{conjectures/A088538/draft.txt}}
{\small\raggedright\ttfamily
\%C A088538 The conjecture of Yalcin Aktar is true: if z = r*exp(i*t) then (z + |z|)/2 = r*cos(t/2)*exp(i*t/2), so u(n) = sqrt(2)*exp(i*Pi/2\char94{}(n+2))*Product\char95{}\char123{}k=1..n\char125{} cos(Pi/2\char94{}(k+2)), which tends to sqrt(2)*sin(Pi/4)/(Pi/4) = 4/Pi by Viète's product. - \char95{}Joaquim Manji-Jung\char95{}, Nov 19 2026\par
}
\subsection*{\texttt{conjectures/A088617/draft.txt}}
{\small\raggedright\ttfamily
\%C A088617 Callan's conjecture is true: the ratio T(n,k+1)/T(n,k) = (n+k+1)(n-k)/((k+1)(k+2)) is decreasing in k and, at k = x*n, tends to (1-x\char94{}2)/x\char94{}2, which equals 1 at x = 1/sqrt(2); hence the distribution of k concentrates around n/sqrt(2), and the expected number of U's is \char126{} n/sqrt(2). - \char95{}Joaquim Manji-Jung\char95{}, Nov 19 2026\par
}
\subsection*{\texttt{conjectures/A112310/draft.txt}}
{\small\raggedright\ttfamily
\%C A112310 Adamson's conjecture is true from the second group on: for k \char62{}= 2 the k-th group consists of the n whose lazy representation has k-1 digits, and its sum is A023610(k-2). Proof: the lazy representations with L digits are the binary words of length L with no 00 and leading digit 1; such a word is 1 followed by a word of length L-1, or 10 followed by a word of length L-2, so the total number T(L) of ones satisfies T(L) = T(L-1) + T(L-2) + A000045(L+1) with T(1) = 1, T(2) = 3, which is the recurrence of A023610. - \char95{}Joaquim Manji-Jung\char95{}, Nov 18 2026\par
}
\subsection*{\texttt{conjectures/A118006/draft.txt}}
{\small\raggedright\ttfamily
\%C A118006 The functional equation of Hebisch and Rubey, as corrected by Somos, is true when f(x) = Sum\char95{}\char123{}n\char62{}=1\char125{} a(n)*x\char94{}(n-1) (it fails with x\char94{}n). Proof: the pairs of letters of the sequence (01 and 10) form the fixed point of a -\char62{} aab, b -\char62{} abb, and the identity a(n) + a(n-1) + a(n-2) - [n = 3m+3] a(m+1) = (2 if n == 4, 5 (mod 6), else 1) for n \char62{}= 2 is then proved with Walnut in base 3. - \char95{}Joaquim Manji-Jung\char95{}, Nov 18 2026\par
}
\subsection*{\texttt{conjectures/A120160/draft.txt}}
{\small\raggedright\ttfamily
\%C A120160 All three conjectures of Hadjicostas are true. The Josephus one is the case k = 5 of the general theorem (see A005427). For Conjectures 1 and 2: in base 5/4 the parent of N \char62{}= 1 is 4*floor(N/5), so numbers with n digits form an interval I\char95{}n, and I\char95{}\char123{}n+1\char125{} is the union of the blocks [5M, 5M+4] over the multiples 4M in I\char95{}n; with U\char95{}n = ceiling(min I\char95{}n/4) one gets U\char95{}\char123{}n+1\char125{} = U\char95{}n + ceiling(U\char95{}n/4) and \#(multiples of 4 in I\char95{}n) = U\char95{}\char123{}n+1\char125{} - U\char95{}n, i.e. U\char95{}n = Sum\char95{}\char123{}i\char60{}=n\char125{} a(i). Conjecture 1 holds when 0 is not counted (it fails only at n = 2 otherwise). - \char95{}Joaquim Manji-Jung\char95{}, Nov 19 2026\par
}
\subsection*{\texttt{conjectures/A155200/draft.txt}}
{\small\raggedright\ttfamily
\%C A155200 The conjecture is true: A(x) = exp(2x)*R(x) with R(x) = exp(Sum\char95{}\char123{}n\char62{}=2\char125{} 2\char94{}(n\char94{}2) x\char94{}n/n). By Legendre, v\char95{}2(2\char94{}j/j!) = A000120(j), and every term of [x\char94{}t] R has 2-adic valuation \char62{}= t+1 (since n\char94{}2 - v\char95{}2(n) \char62{}= n+1 for n \char62{}= 2 and v\char95{}2(k!) \char60{}= k-1); hence in a(n) = Sum\char95{}t 2\char94{}(n-t)/(n-t)! * r(t) the term t = 0 has valuation A000120(n) and all others have larger valuation. - \char95{}Joaquim Manji-Jung\char95{}, Nov 19 2026\par
}
\subsection*{\texttt{conjectures/A161580/draft.txt}}
{\small\raggedright\ttfamily
\%C A161580 The conjecture is true. It is a first-order statement in (k, n) over the Thue-Morse sequence, and the Walnut theorem prover confirms it for all odd k at once, both for the positions with A010060(n) + A010060(n+k) = 1 and for those with A010060(n) = A010060(n+k); for the latter see also Allouche (2015). - \char95{}Joaquim Manji-Jung\char95{}, Nov 18 2026\par
\%H A161580 J.-P. Allouche, \char60{}a href="https://doi.org/10.5802/jtnb.906"\char62{}Thue, Combinatorics on words, and conjectures inspired by the Thue-Morse sequence\char60{}/a\char62{}, J. de Théorie des Nombres de Bordeaux, 27, no. 2 (2015), 375-388.\par
}
\subsection*{\texttt{conjectures/A170941/draft.txt}}
{\small\raggedright\ttfamily
\%C A170941 Callan's conjecture is true: by inclusion-exclusion a(n) = Sum\char95{}h (-1)\char94{}h*binomial(n-h,h)*A000085(n-2h), and since A000085(m) \char62{}= (m-1)*A000085(m-2) each term divided by A000085(n) is at most 1/h!, while it tends to 1/h! for fixed h (A000085(m)/A000085(m-1) \char126{} sqrt(m)); dominated convergence gives the limit 1/e. - \char95{}Joaquim Manji-Jung\char95{}, Nov 19 2026\par
}
\subsection*{\texttt{conjectures/A225132/draft.txt}}
{\small\raggedright\ttfamily
\%C A225132 The conjecture is true. The data are the coefficients of x/g(x) = 1/h(x), h(x) = Sum\char95{}\char123{}n\char62{}=0\char125{} A010060(n+1)*x\char94{}n (not of 1/(1+g(x))). A certified argument-principle computation shows that h has exactly one zero in |x| \char60{} 0.76, the simple real zero x0 = -0.71165406530823624..., so a(n) \char126{} -x0\char94{}(-n-1)/h'(x0) and a(n+1)/a(n) -\char62{} 1/x0 = -1.40517710605204438... - \char95{}Joaquim Manji-Jung\char95{}, Nov 19 2026\par
}
\subsection*{\texttt{conjectures/A226898/draft.txt}}
{\small\raggedright\ttfamily
\%C A226898 Greathouse's conjecture is true. If an odd prime q divides n, split the divisors into the chains d, d*q, ..., d*q\char94{}v (v the q-adic valuation of n); each chain has at least 2 elements, and an interval [u, e*u] contains at most one element of each chain since q \char62{}= 3 \char62{} e, so a(n) \char60{}= d(n)/2. If n = 2\char94{}v, then a(n) \char60{}= 2 \char60{}= (v+1)/2 for v \char62{}= 3. - \char95{}Joaquim Manji-Jung\char95{}, Nov 19 2026\par
}
\subsection*{\texttt{conjectures/A285126/draft.txt}}
{\small\raggedright\ttfamily
\%C A285126 Kimberling's conjecture is false: a(3) = 3 and 3*r - 3 = 1.1711... \char62{} 1. In fact n*r - a(n) is unbounded, since the morphism 0-\char62{}10, 1-\char62{}0000 is not Pisot (its second eigenvalue (1-sqrt(17))/2 has absolute value \char62{} 1); n*r - a(n) exceeds 1000 for some n \char60{} 2.2*10\char94{}6. - \char95{}Joaquim Manji-Jung\char95{}, Nov 19 2026\par
\%C A285127 Kimberling's conjecture is false: a(2) = 11 and 2*r - 11 = -3.8768... \char60{} -3; moreover n*r - a(n) is unbounded (the morphism 0-\char62{}10, 1-\char62{}0000 is not Pisot). - \char95{}Joaquim Manji-Jung\char95{}, Nov 19 2026\par
}
\subsection*{\texttt{conjectures/A285503/draft.txt}}
{\small\raggedright\ttfamily
\%C A285503 The conjecture is true: the underlying morphism is primitive with incidence matrix [[0,2],[2,2]], whose Perron eigenvector (2, 1+sqrt(5)) gives the frequency (sqrt(5)-1)/2 = 1/phi for the letter 1, so a(n)/n -\char62{} phi. - \char95{}Joaquim Manji-Jung\char95{}, Nov 19 2026\par
\%C A285517 The conjecture is true: the underlying morphism is primitive with incidence matrix [[0,2],[2,2]], whose Perron eigenvector (2, 1+sqrt(5)) gives the frequency (sqrt(5)-1)/2 = 1/phi for the letter 1, so a(n)/n -\char62{} phi. - \char95{}Joaquim Manji-Jung\char95{}, Nov 19 2026\par
\%C A285567 The conjecture is true: the underlying morphism is primitive with incidence matrix [[0,2],[2,2]], whose Perron eigenvector (2, 1+sqrt(5)) gives the frequency (sqrt(5)-1)/2 = 1/phi for the letter 1, so a(n)/n -\char62{} phi. - \char95{}Joaquim Manji-Jung\char95{}, Nov 19 2026\par
\%C A285623 The conjecture is true: the underlying morphism is primitive with incidence matrix [[0,2],[2,2]], whose Perron eigenvector (2, 1+sqrt(5)) gives the frequency (sqrt(5)-1)/2 = 1/phi for the letter 1, so a(n)/n -\char62{} phi. - \char95{}Joaquim Manji-Jung\char95{}, Nov 19 2026\par
\%C A285627 The conjecture is true: the underlying morphism is primitive with incidence matrix [[0,2],[2,2]], whose Perron eigenvector (2, 1+sqrt(5)) gives the frequency (sqrt(5)-1)/2 = 1/phi for the letter 1, so a(n)/n -\char62{} phi. - \char95{}Joaquim Manji-Jung\char95{}, Nov 19 2026\par
\%C A285663 The conjecture is true: the underlying morphism is primitive with incidence matrix [[0,2],[2,2]], whose Perron eigenvector (2, 1+sqrt(5)) gives the frequency (sqrt(5)-1)/2 = 1/phi for the letter 1, so a(n)/n -\char62{} phi. - \char95{}Joaquim Manji-Jung\char95{}, Nov 19 2026\par
}
\subsection*{\texttt{conjectures/A297464/draft.txt}}
{\small\raggedright\ttfamily
\%C A297464 The conjectured inequalities hold when n is replaced by n+1 (as written, with offset 0, they fail at n = 0), and the conjectured recurrence a(n) = a(n-1) + a(n-3) - a(n-4) is true: a(n) - floor(10n/3) and the analogous differences for b, c, d are 3-automatic, and Walnut verifies that these candidates partition the positive integers and obey the "least new" rules, hence are the solution. - \char95{}Joaquim Manji-Jung\char95{}, Nov 19 2026\par
\%C A297465 See A297464: the conjectured inequalities are true after replacing n by n+1 (Walnut proof); as written (offset 0) they fail at n = 0. - \char95{}Joaquim Manji-Jung\char95{}, Nov 19 2026\par
\%C A297466 See A297464: the conjectured inequalities are true after replacing n by n+1 (Walnut proof); as written (offset 0) they fail at n = 0. - \char95{}Joaquim Manji-Jung\char95{}, Nov 19 2026\par
}
\subsection*{\texttt{conjectures/A297469/draft.txt}}
{\small\raggedright\ttfamily
\%C A297469 Sloane's conjecture is true, and so are (2), (4) and cc = A047218 (shifted): the three sequences given by Sloane's formula (bb), cc(2k) = 5k+3, cc(2k+1) = 5k+5, and aa(n) = bb(n) - cc(n-1) are 2-automatic, and Walnut verifies that they partition the positive integers and satisfy the "least new" rules, so by induction they are the solution. Statements (1) and (3) are false as written (n = 1 gives -4 and -5); the true forms are 2*aa(n) - 5*n in \char123{}2,3,4\char125{} and bb(n) - 5*n in \char123{}1,2\char125{}. - \char95{}Joaquim Manji-Jung\char95{}, Nov 19 2026\par
\%C A298468 Sloane's conjecture is true (Walnut proof via 2-automatic candidates and the "least new" rules; see A297469). - \char95{}Joaquim Manji-Jung\char95{}, Nov 19 2026\par
}
\subsection*{\texttt{conjectures/A298173/draft.txt}}
{\small\raggedright\ttfamily
\%C A298173 Both observations are true: a(n) - 3*a(n-1) = Sum\char95{}\char123{}j=0..n\char125{} (-1)\char94{}j b(n-j) lies between ceiling(n/2) (pair the terms; b is increasing) and b(n) \char60{}= 2n+1 (only 2n numbers are placed before b(n)), so it is unbounded and a(n)/a(n-1) = 3 + O(n/3\char94{}n). - \char95{}Joaquim Manji-Jung\char95{}, Nov 19 2026\par
\%C A298741 Both observations are true: a(n) - 2*a(n-1) = b(n) + Sum\char95{}\char123{}j=0..n\char125{} (-1)\char94{}j b(n-j) lies between n+1 and 4n+2 (b is increasing and b(n) \char60{}= 2n+1), so it is unbounded and a(n)/a(n-1) -\char62{} 2. - \char95{}Joaquim Manji-Jung\char95{}, Nov 19 2026\par
}
\subsection*{\texttt{conjectures/A299405/draft.txt}}
{\small\raggedright\ttfamily
\%C A299405 The conjecture is false: even with n replaced by n+1 (as written, with offset 0, every line fails at n = 0), c(26) = 113 gives 17*27 - 3 - 4*113 = 4 \char62{} 3. With that shift, the bounds for a, b, d, e are true, and the bound for c is true with 4 in place of 3 (Walnut proof: the differences x(n) - floor(17n/4) are 2-automatic and the candidates obey the "least new" rules). - \char95{}Joaquim Manji-Jung\char95{}, Nov 19 2026\par
\%C A299409 The conjecture is false: c(26) = 113 violates the bound for c even after shifting n to n+1; see A299405 (Walnut proof of the corrected bounds). - \char95{}Joaquim Manji-Jung\char95{}, Nov 19 2026\par
\%C A299637 The conjecture is false: c(26) = 113 violates the bound for c even after shifting n to n+1; see A299405 (Walnut proof of the corrected bounds). - \char95{}Joaquim Manji-Jung\char95{}, Nov 19 2026\par
\%C A299638 The conjecture is false: c(26) = 113 violates the bound for c even after shifting n to n+1; see A299405 (Walnut proof of the corrected bounds). - \char95{}Joaquim Manji-Jung\char95{}, Nov 19 2026\par
\%C A299641 The conjecture is false: c(26) = 113 violates the bound for c even after shifting n to n+1; see A299405 (Walnut proof of the corrected bounds). - \char95{}Joaquim Manji-Jung\char95{}, Nov 19 2026\par
}
\subsection*{\texttt{conjectures/A304451/draft.txt}}
{\small\raggedright\ttfamily
\%C A304451 The limits a(n)/n -\char62{} 5/2 and b(n)/n -\char62{} 5/3 are true: b(n) = n + 1 + \#\char123{}k: a(k) \char60{} b(n)\char125{}, so with beta+ = limsup b(n)/n and beta- = liminf b(n)/n one gets beta+ \char60{}= f(beta-) and beta- \char62{}= f(beta+) for f(x) = 3x/(3x-2), which forces beta+ = beta- = 5/3. The inequalities fail at n = 1 (a(1) = 4, b(1) = 3) and n = 2 (a(2) = 8), and hold for n \char62{}= 3 (for all n with n+1 in place of n). - \char95{}Joaquim Manji-Jung\char95{}, Nov 19 2026\par
\%C A304452 See A304451: the limits are true; the inequality for b fails at n = 1 (b(1) = 3) and holds for n \char62{}= 2. - \char95{}Joaquim Manji-Jung\char95{}, Nov 19 2026\par
}
\subsection*{\texttt{conjectures/A347171/draft.txt}}
{\small\raggedright\ttfamily
\%C A347171 The conjecture of Doche and Mendès France is true: for every n, P\char95{}n has exactly A285869(n) real roots, all simple. For y \char60{} 0 write t = sqrt(-2y) and x = (1-y)/t = 2*cos(theta); then P\char95{}n(y) = t\char94{}n*(S(n,x) - t*S(n-1,x)) with Chebyshev S, and the zeros are the solutions of n*theta + arg(exp(i*theta) - t) = k*Pi on the two branches y \char60{} -1 and -1 \char60{} y \char60{} -3+2*sqrt(2), where this phase is respectively strictly concave and strictly increasing. There are no zeros with x \char62{} 2, and for y \char62{} 0 exactly one zero if n is even and none if n is odd. See the link. - \char95{}Joaquim Manji-Jung\char95{}, Nov 18 2026\par
}
\subsection*{\texttt{conjectures/A357097/draft.txt}}
{\small\raggedright\ttfamily
\%C A357097 The conjectured formula A(A019586(x), A019586(y)) = A019586(A101330(x,y)) is true. Proof: with w = phi - 1, the product of pairs is the multiplication of j + k*w in Z[phi], and S\char95{}n = \char123{}b\char95{}n*phi\char94{}m\char125{}. For n \char62{}= 1 with Zeckendorf representation Sum F(c\char95{}i), the pair (successor of n in its Wythoff row, n) corresponds to t(n) = Sum phi\char94{}(c\char95{}i) = Sum F(c\char95{}i - 1) + n*phi. The conjugate of t(n) lies in the interval (-1/phi\char94{}2, 1/phi) of length 1, which determines t(n) among the elements a + n*phi, and the product of two numbers of this interval is again in it; since the coefficient of phi in t(x)*t(y) is the Knuth product A101330(x,y), t(x)*t(y) = t(A101330(x,y)). - \char95{}Joaquim Manji-Jung\char95{}, Nov 18 2026\par
}
\subsection*{\texttt{conjectures/A358360/draft.txt}}
{\small\raggedright\ttfamily
\%C A358360 The five conjectures of Oudra are true; they follow from Theorems 4 and 6 of Alekseyev, Amdeberhan, Shallit and Vukusic. Theorem 4 gives a(2n) = A000989(n) and a(2n+1) = A000989(n) + 1 + A007949(2n+1), and A254046(n+1) = 1 + A007949(2n+1); the formula with A168570 follows from A168570(2j) = 1 + A007949(j) and binomial(2m+2,m+1) = binomial(2m,m)*2*(2m+1)/(m+1). Theorem 6 gives a(3n+r) = a(n) + ((3n+r) mod 2), which is the recurrence of Sum\char95{}\char123{}k\char62{}=0\char125{} (floor(n/3\char94{}k) mod 2); and the alternating sum of A051064 equals that sum after exchanging the order of summation. - \char95{}Joaquim Manji-Jung\char95{}, Nov 18 2026\par
}
\subsection*{\texttt{conjectures/A361190/draft.txt}}
{\small\raggedright\ttfamily
\%C A361190 Kotesovec's conjectured exponent follows from the theorem of Denisov and Wachtel on random walks in cones: in the coordinates (x-y, y) the steps become (0,1), (2,-1), (-1,0) in the quarter plane; the exponential change of measure with weights 1/4, 1/4, 1/2 removes the drift (growth rate 2*sqrt(2) per step, i.e., 64\char94{}n), and the resulting walk has correlation coefficient -1/sqrt(3), so the cone angle is arccos(1/sqrt(3)) = arctan(sqrt(2)) and the excursion exponent is 1 + Pi/arctan(sqrt(2)). Since this is irrational, the g.f. is not D-finite. - \char95{}Joaquim Manji-Jung\char95{}, Nov 19 2026\par
}
\subsection*{\texttt{conjectures/A363308/draft.txt}}
{\small\raggedright\ttfamily
\%C A363308 The conjecture is true. Modulo 2, C(x) == P(x)/x where P(x) = Sum\char95{}\char123{}k\char62{}=0\char125{} x\char94{}(2\char94{}k), and P\char94{}2 = P + x. With y = x*C(x)\char94{}3 and A = 1 + P one checks y*A\char94{}2 = A - 1 over GF(2); since C(y) is the unique power series with F = 1 + y*F\char94{}2, the g.f. is congruent to 1 + Sum\char95{}\char123{}k\char62{}=0\char125{} x\char94{}(2\char94{}k) (mod 2). - \char95{}Joaquim Manji-Jung\char95{}, Nov 19 2026\par
}
\subsection*{\texttt{conjectures/A366860/draft.txt}}
{\small\raggedright\ttfamily
\%C A366860 The conjecture is true: the first differences a(n) - a(n-1), n \char62{}= 1, form the fixed point of the morphism 0 -\char62{} 1\char94{}3, 1 -\char62{} 1\char94{}30 starting with 1 (each block of the fixed point begins with 3 ones and ends with a 0 exactly when its preimage letter is 1, which gives the recurrence directly). The same holds for a(n) = n - a(floor(a(n-1)/k)) for every k \char62{}= 1, so a(n)/n -\char62{} (sqrt(k\char94{}2+4k) - k)/2. See the link for a proof. - \char95{}Joaquim Manji-Jung\char95{}, Nov 18 2026\par
\par
\%C A366870 The conjecture is true: the first differences a(n) - a(n-1), n \char62{}= 1, form the fixed point of the morphism 0 -\char62{} 1\char94{}4, 1 -\char62{} 1\char94{}40 starting with 1 (each block of the fixed point begins with 4 ones and ends with a 0 exactly when its preimage letter is 1, which gives the recurrence directly). The same holds for a(n) = n - a(floor(a(n-1)/k)) for every k \char62{}= 1, so a(n)/n -\char62{} (sqrt(k\char94{}2+4k) - k)/2. See the link for a proof. - \char95{}Joaquim Manji-Jung\char95{}, Nov 18 2026\par
\par
\%C A366871 The conjecture is true: the first differences a(n) - a(n-1), n \char62{}= 1, form the fixed point of the morphism 0 -\char62{} 1\char94{}5, 1 -\char62{} 1\char94{}50 starting with 1 (each block of the fixed point begins with 5 ones and ends with a 0 exactly when its preimage letter is 1, which gives the recurrence directly). The same holds for a(n) = n - a(floor(a(n-1)/k)) for every k \char62{}= 1, so a(n)/n -\char62{} (sqrt(k\char94{}2+4k) - k)/2. See the link for a proof. - \char95{}Joaquim Manji-Jung\char95{}, Nov 18 2026\par
}
\subsection*{\texttt{conjectures/A368281/draft.txt}}
{\small\raggedright\ttfamily
\%C A368281 The conjecture is true for all n: the fixed point w of 0 -\char62{} 1\char94{}n, 1 -\char62{} 1\char94{}n 0 is the concatenation of the blocks 1\char94{}n and 1\char94{}n 0, so if the prefix of length k-1 consists of j complete blocks and r letters, its sum is nj + r with 0 \char60{}= r \char60{}= n, and the number of 0's in the prefix of length k equals the partial sum up to floor(b(k-1)/n), which is the recurrence of A368282. See A366860 for details. - \char95{}Joaquim Manji-Jung\char95{}, Nov 19 2026\par
}
\subsection*{\texttt{conjectures/A371176/draft.txt}}
{\small\raggedright\ttfamily
\%C A371176 Kotesovec's conjecture is true, and in fact lim\char95{}\char123{}n-\char62{}oo\char125{} log(a(n))/log(n) = log(2)/log((1+sqrt(5))/2). Proof: by Branicky's correspondence with Schreier sets, the number of terms in [2\char94{}(m-1), 2\char94{}m) is the number of Schreier sets with maximum m, which is A000045(m); hence there are A000045(m+2) - 1 terms \char60{} 2\char94{}m, so if a(n) has m binary digits then log(n) = m*log((1+sqrt(5))/2) + O(1) and log(a(n)) = m*log(2) + O(1). - \char95{}Joaquim Manji-Jung\char95{}, Nov 18 2026\par
}
\subsection*{\texttt{conjectures/batch8/draft.txt}}
{\small\raggedright\ttfamily
\%C A373308 The conjecture is true: Product\char95{}\char123{}n\char62{}=0\char125{} (1 - x\char94{}(2\char94{}n)) = Sum\char95{}n (-1)\char94{}A010060(n) x\char94{}n, and cubing over GF(3) gives Sum\char95{}n (-1)\char94{}A010060(n) x\char94{}(3n), so a(3n) == (-1)\char94{}A010060(n) == A001285(n) (mod 3) and a(n) == 0 (mod 3) if 3 does not divide n. - \char95{}Joaquim Manji-Jung\char95{}, Nov 19 2026\par
\%C A000918 Jeffery's conjecture is true: A007814(A000108(m)) = A000120(m+1) - 1 (Kummer), and the least m+1 with binary weight n is 2\char94{}n - 1. - \char95{}Joaquim Manji-Jung\char95{}, Nov 19 2026\par
\%C A048881 Kimberling's conjecture is true: the n-th word of A076478 is the binary expansion of n+1 with its leading 1 removed. - \char95{}Joaquim Manji-Jung\char95{}, Nov 19 2026\par
\%C A059906 Yanev's conjecture is true: in the second formula the terms for k = 2m-1 and k = 2m cancel except for the contribution 2\char94{}((i-1)/2) of each bit of n in odd position i, and the first formula follows from it by summation by parts. - \char95{}Joaquim Manji-Jung\char95{}, Nov 19 2026\par
\%C A073089 Yanev's conjecture is true: A057661(m) = (1 + Sum\char95{}\char123{}d|m\char125{} d*phi(d))/2, and Sum\char95{}\char123{}d|m\char125{} d*phi(d) == 3\char94{}([m even] + [odd part of m == 3 (mod 4)]) (mod 4), which matches a(n) = A014577(n-2) + A056594(n). - \char95{}Joaquim Manji-Jung\char95{}, Nov 19 2026\par
\%C A284388 George Beck's conjecture is true: with sigma\char94{}2: 0 -\char62{} 001, 1 -\char62{} 11001 and rho: 0 -\char62{} 010, 1 -\char62{} 01110 (the run-length morphism of Thue-Morse, see A026465, with lengths decreased by 1), one has 01*sigma\char94{}2(a) = rho(a)*01 for each letter, so writing this word as 0Y gives Y = 01*sigma\char94{}2(Y) = rho(Y), the fixed point of rho. - \char95{}Joaquim Manji-Jung\char95{}, Nov 19 2026\par
}
\subsection*{\texttt{conjectures/beattyeq/draft.txt}}
{\small\raggedright\ttfamily
\%C A297836 The conjecture is false: a(2) = 11 and 11 - (5 + sqrt(13)) = 2.394... \char62{} 2; there are 605551 violations for n \char60{} 2*10\char94{}6 (max 2.8685...). - \char95{}Joaquim Manji-Jung\char95{}, Nov 19 2026\par
\%C A298112 The conjecture as written is false: b(n)/n -\char62{} (-1 + sqrt(13))/2 (forced by 1/alpha + 1/beta = 1 with alpha = beta + 3), not sqrt(2), so b(n) - n*sqrt(2) -\char62{} -infinity (already b(7) = 10 gives 0.10 \char60{} 2/5). With L = (-1 + sqrt(13))/2 in place of sqrt(2), 2/5 \char60{} b(n) - n*L \char60{} 3 holds for n \char60{} 2*10\char94{}6. - \char95{}Joaquim Manji-Jung\char95{}, Nov 19 2026\par
\%C A298113 The conjecture as written is false: b(n)/n -\char62{} -1 + sqrt(5), not sqrt(2), so b(n) - n*sqrt(2) -\char62{} -infinity (b(7) = 10 gives 0.10 \char60{} 1/2). With L = -1 + sqrt(5), 1/2 \char60{} b(n) - n*L \char60{} 3 holds for n \char60{} 2*10\char94{}6. - \char95{}Joaquim Manji-Jung\char95{}, Nov 19 2026\par
\%C A297834 The data do not satisfy the definition from a(3) on: with b(1) = 4 and b(2) = 6 (A298109), a(3) = 2*6 - 4 + 2*3 - 4 = 10, not 8. The solution of the titular equation begins 1, 2, 5, 10, 12, 15, 18, 23, 27, 29, 34, 36, ... - \char95{}Joaquim Manji-Jung\char95{}, Nov 19 2026\par
}
\subsection*{\texttt{conjectures/compeq/draft.txt}}
{\small\raggedright\ttfamily
\%C A293057 The conjecture is true: the forcing term g(n) = a(n) - a(n-1) - a(n-2) is \char62{}= 0, so a grows exponentially and b(n) \char60{}= 2n+1 for large n; then e(n) = a(n) - phi*a(n-1) satisfies e(n) = -e(n-1)/phi + g(n) and is polynomially bounded since g(n) is, so a(n)/a(n-1) - phi = e(n)/a(n-1) -\char62{} 0. - \char95{}Joaquim Manji-Jung\char95{}, Nov 19 2026\par
\%C A293058 The conjecture is true: the forcing term g(n) = a(n) - a(n-1) - a(n-2) is \char62{}= 0, so a grows exponentially and b(n) \char60{}= 2n+1 for large n; then e(n) = a(n) - phi*a(n-1) satisfies e(n) = -e(n-1)/phi + g(n) and is polynomially bounded since g(n) is, so a(n)/a(n-1) - phi = e(n)/a(n-1) -\char62{} 0. - \char95{}Joaquim Manji-Jung\char95{}, Nov 19 2026\par
\%C A293076 The conjecture is true: the forcing term g(n) = a(n) - a(n-1) - a(n-2) is \char62{}= 0, so a grows exponentially and b(n) \char60{}= 2n+1 for large n; then e(n) = a(n) - phi*a(n-1) satisfies e(n) = -e(n-1)/phi + g(n) and is polynomially bounded since g(n) is, so a(n)/a(n-1) - phi = e(n)/a(n-1) -\char62{} 0. - \char95{}Joaquim Manji-Jung\char95{}, Nov 19 2026\par
\%C A293316 The conjecture is true: the forcing term g(n) = a(n) - a(n-1) - a(n-2) is \char62{}= 0, so a grows exponentially and b(n) \char60{}= 2n+1 for large n; then e(n) = a(n) - phi*a(n-1) satisfies e(n) = -e(n-1)/phi + g(n) and is polynomially bounded since g(n) is, so a(n)/a(n-1) - phi = e(n)/a(n-1) -\char62{} 0. - \char95{}Joaquim Manji-Jung\char95{}, Nov 19 2026\par
\%C A293317 The conjecture is true: the forcing term g(n) = a(n) - a(n-1) - a(n-2) is \char62{}= 0, so a grows exponentially and b(n) \char60{}= 2n+1 for large n; then e(n) = a(n) - phi*a(n-1) satisfies e(n) = -e(n-1)/phi + g(n) and is polynomially bounded since g(n) is, so a(n)/a(n-1) - phi = e(n)/a(n-1) -\char62{} 0. - \char95{}Joaquim Manji-Jung\char95{}, Nov 19 2026\par
\%C A293349 The conjecture is true: the forcing term g(n) = a(n) - a(n-1) - a(n-2) is \char62{}= 0, so a grows exponentially and b(n) \char60{}= 2n+1 for large n; then e(n) = a(n) - phi*a(n-1) satisfies e(n) = -e(n-1)/phi + g(n) and is polynomially bounded since g(n) is, so a(n)/a(n-1) - phi = e(n)/a(n-1) -\char62{} 0. - \char95{}Joaquim Manji-Jung\char95{}, Nov 19 2026\par
\%C A293350 The conjecture is true: the forcing term g(n) = a(n) - a(n-1) - a(n-2) is \char62{}= 0, so a grows exponentially and b(n) \char60{}= 2n+1 for large n; then e(n) = a(n) - phi*a(n-1) satisfies e(n) = -e(n-1)/phi + g(n) and is polynomially bounded since g(n) is, so a(n)/a(n-1) - phi = e(n)/a(n-1) -\char62{} 0. - \char95{}Joaquim Manji-Jung\char95{}, Nov 19 2026\par
\%C A293351 The conjecture is true: the forcing term g(n) = a(n) - a(n-1) - a(n-2) is \char62{}= 0, so a grows exponentially and b(n) \char60{}= 2n+1 for large n; then e(n) = a(n) - phi*a(n-1) satisfies e(n) = -e(n-1)/phi + g(n) and is polynomially bounded since g(n) is, so a(n)/a(n-1) - phi = e(n)/a(n-1) -\char62{} 0. - \char95{}Joaquim Manji-Jung\char95{}, Nov 19 2026\par
\%C A293357 The conjecture is true: the forcing term g(n) = a(n) - a(n-1) - a(n-2) is \char62{}= 0, so a grows exponentially and b(n) \char60{}= 2n+1 for large n; then e(n) = a(n) - phi*a(n-1) satisfies e(n) = -e(n-1)/phi + g(n) and is polynomially bounded since g(n) is, so a(n)/a(n-1) - phi = e(n)/a(n-1) -\char62{} 0. - \char95{}Joaquim Manji-Jung\char95{}, Nov 19 2026\par
\%C A293358 The conjecture is true: the forcing term g(n) = a(n) - a(n-1) - a(n-2) is \char62{}= 0, so a grows exponentially and b(n) \char60{}= 2n+1 for large n; then e(n) = a(n) - phi*a(n-1) satisfies e(n) = -e(n-1)/phi + g(n) and is polynomially bounded since g(n) is, so a(n)/a(n-1) - phi = e(n)/a(n-1) -\char62{} 0. - \char95{}Joaquim Manji-Jung\char95{}, Nov 19 2026\par
\%C A293406 The conjecture is true: the forcing term g(n) = a(n) - a(n-1) - a(n-2) is \char62{}= 0, so a grows exponentially and b(n) \char60{}= 2n+1 for large n; then e(n) = a(n) - phi*a(n-1) satisfies e(n) = -e(n-1)/phi + g(n) and is polynomially bounded since g(n) is, so a(n)/a(n-1) - phi = e(n)/a(n-1) -\char62{} 0. - \char95{}Joaquim Manji-Jung\char95{}, Nov 19 2026\par
\%C A293765 The conjecture is true: the forcing term g(n) = a(n) - a(n-1) - a(n-2) is \char62{}= 0, so a grows exponentially and b(n) \char60{}= 2n+1 for large n; then e(n) = a(n) - phi*a(n-1) satisfies e(n) = -e(n-1)/phi + g(n) and is polynomially bounded since g(n) is, so a(n)/a(n-1) - phi = e(n)/a(n-1) -\char62{} 0. - \char95{}Joaquim Manji-Jung\char95{}, Nov 19 2026\par
\%C A293766 The conjecture is true: the forcing term g(n) = a(n) - a(n-1) - a(n-2) is \char62{}= 0, so a grows exponentially and b(n) \char60{}= 2n+1 for large n; then e(n) = a(n) - phi*a(n-1) satisfies e(n) = -e(n-1)/phi + g(n) and is polynomially bounded since g(n) is, so a(n)/a(n-1) - phi = e(n)/a(n-1) -\char62{} 0. - \char95{}Joaquim Manji-Jung\char95{}, Nov 19 2026\par
\%C A293767 The conjecture is true: the forcing term g(n) = a(n) - a(n-1) - a(n-2) is \char62{}= 0, so a grows exponentially and b(n) \char60{}= 2n+1 for large n; then e(n) = a(n) - phi*a(n-1) satisfies e(n) = -e(n-1)/phi + g(n) and is polynomially bounded since g(n) is, so a(n)/a(n-1) - phi = e(n)/a(n-1) -\char62{} 0. - \char95{}Joaquim Manji-Jung\char95{}, Nov 19 2026\par
\%C A294170 Proof of the limit: the forcing term g(n) = a(n) - a(n-1) - a(n-2) is \char62{}= 0, so a grows exponentially and b(n) \char60{}= 2n+1 for large n; then e(n) = a(n) - phi*a(n-1) satisfies e(n) = -e(n-1)/phi + g(n) and is polynomially bounded since g(n) is, so a(n)/a(n-1) - phi = e(n)/a(n-1) -\char62{} 0. - \char95{}Joaquim Manji-Jung\char95{}, Nov 19 2026\par
\%C A294365 The conjecture is true: the forcing term g(n) = a(n) - a(n-1) - a(n-2) is \char62{}= 0, so a grows exponentially and b(n) \char60{}= 2n+1 for large n; then e(n) = a(n) - phi*a(n-1) satisfies e(n) = -e(n-1)/phi + g(n) and is polynomially bounded since g(n) is, so a(n)/a(n-1) - phi = e(n)/a(n-1) -\char62{} 0. - \char95{}Joaquim Manji-Jung\char95{}, Nov 19 2026\par
\%C A294366 The conjecture is true: the forcing term g(n) = a(n) - a(n-1) - a(n-2) is \char62{}= 0, so a grows exponentially and b(n) \char60{}= 2n+1 for large n; then e(n) = a(n) - phi*a(n-1) satisfies e(n) = -e(n-1)/phi + g(n) and is polynomially bounded since g(n) is, so a(n)/a(n-1) - phi = e(n)/a(n-1) -\char62{} 0. - \char95{}Joaquim Manji-Jung\char95{}, Nov 19 2026\par
\%C A294367 The conjecture is true: the forcing term g(n) = a(n) - a(n-1) - a(n-2) is \char62{}= 0, so a grows exponentially and b(n) \char60{}= 2n+1 for large n; then e(n) = a(n) - phi*a(n-1) satisfies e(n) = -e(n-1)/phi + g(n) and is polynomially bounded since g(n) is, so a(n)/a(n-1) - phi = e(n)/a(n-1) -\char62{} 0. - \char95{}Joaquim Manji-Jung\char95{}, Nov 19 2026\par
\%C A294368 The conjecture is true: the forcing term g(n) = a(n) - a(n-1) - a(n-2) is \char62{}= 0, so a grows exponentially and b(n) \char60{}= 2n+1 for large n; then e(n) = a(n) - phi*a(n-1) satisfies e(n) = -e(n-1)/phi + g(n) and is polynomially bounded since g(n) is, so a(n)/a(n-1) - phi = e(n)/a(n-1) -\char62{} 0. - \char95{}Joaquim Manji-Jung\char95{}, Nov 19 2026\par
\%C A294413 The conjecture is true: by induction a(n) \char62{}= 4n for n \char62{}= 5, hence b(n-1) \char60{}= 2n-1 and a(n) \char62{}= a(n-1) + a(n-2) - 2n, so a(n) - 2n - 6 grows like phi\char94{}n; then e(n) = a(n) - phi*a(n-1) satisfies e(n) = -e(n-1)/phi + g(n) and is polynomially bounded since g(n) is, so a(n)/a(n-1) - phi = e(n)/a(n-1) -\char62{} 0. - \char95{}Joaquim Manji-Jung\char95{}, Nov 19 2026\par
\%C A294414 The conjecture is true: the forcing term g(n) = a(n) - a(n-1) - a(n-2) is \char62{}= 0, so a grows exponentially and b(n) \char60{}= 2n+1 for large n; then e(n) = a(n) - phi*a(n-1) satisfies e(n) = -e(n-1)/phi + g(n) and is polynomially bounded since g(n) is, so a(n)/a(n-1) - phi = e(n)/a(n-1) -\char62{} 0. - \char95{}Joaquim Manji-Jung\char95{}, Nov 19 2026\par
\%C A294415 The conjecture is true: the forcing term g(n) = a(n) - a(n-1) - a(n-2) is \char62{}= 0, so a grows exponentially and b(n) \char60{}= 2n+1 for large n; then e(n) = a(n) - phi*a(n-1) satisfies e(n) = -e(n-1)/phi + g(n) and is polynomially bounded since g(n) is, so a(n)/a(n-1) - phi = e(n)/a(n-1) -\char62{} 0. - \char95{}Joaquim Manji-Jung\char95{}, Nov 19 2026\par
\%C A294416 The conjecture is true: the forcing term g(n) = a(n) - a(n-1) - a(n-2) is \char62{}= 0, so a grows exponentially and b(n) \char60{}= 2n+1 for large n; then e(n) = a(n) - phi*a(n-1) satisfies e(n) = -e(n-1)/phi + g(n) and is polynomially bounded since g(n) is, so a(n)/a(n-1) - phi = e(n)/a(n-1) -\char62{} 0. - \char95{}Joaquim Manji-Jung\char95{}, Nov 19 2026\par
\%C A294417 The conjecture is true: the forcing term g(n) = a(n) - a(n-1) - a(n-2) is \char62{}= 0, so a grows exponentially and b(n) \char60{}= 2n+1 for large n; then e(n) = a(n) - phi*a(n-1) satisfies e(n) = -e(n-1)/phi + g(n) and is polynomially bounded since g(n) is, so a(n)/a(n-1) - phi = e(n)/a(n-1) -\char62{} 0. - \char95{}Joaquim Manji-Jung\char95{}, Nov 19 2026\par
\%C A294418 The conjecture is true: the forcing term g(n) = a(n) - a(n-1) - a(n-2) is \char62{}= 0, so a grows exponentially and b(n) \char60{}= 2n+1 for large n; then e(n) = a(n) - phi*a(n-1) satisfies e(n) = -e(n-1)/phi + g(n) and is polynomially bounded since g(n) is, so a(n)/a(n-1) - phi = e(n)/a(n-1) -\char62{} 0. - \char95{}Joaquim Manji-Jung\char95{}, Nov 19 2026\par
\%C A294419 The conjecture is true: the forcing term g(n) = a(n) - a(n-1) - a(n-2) is \char62{}= 0, so a grows exponentially and b(n) \char60{}= 2n+1 for large n; then e(n) = a(n) - phi*a(n-1) satisfies e(n) = -e(n-1)/phi + g(n) and is polynomially bounded since g(n) is, so a(n)/a(n-1) - phi = e(n)/a(n-1) -\char62{} 0. - \char95{}Joaquim Manji-Jung\char95{}, Nov 19 2026\par
\%C A294420 The conjecture is true: the forcing term g(n) = a(n) - a(n-1) - a(n-2) is \char62{}= 0, so a grows exponentially and b(n) \char60{}= 2n+1 for large n; then e(n) = a(n) - phi*a(n-1) satisfies e(n) = -e(n-1)/phi + g(n) and is polynomially bounded since g(n) is, so a(n)/a(n-1) - phi = e(n)/a(n-1) -\char62{} 0. - \char95{}Joaquim Manji-Jung\char95{}, Nov 19 2026\par
\%C A294421 The conjecture is true: the forcing term g(n) = a(n) - a(n-1) - a(n-2) is \char62{}= 0, so a grows exponentially and b(n) \char60{}= 2n+1 for large n; then e(n) = a(n) - phi*a(n-1) satisfies e(n) = -e(n-1)/phi + g(n) and is polynomially bounded since g(n) is, so a(n)/a(n-1) - phi = e(n)/a(n-1) -\char62{} 0. - \char95{}Joaquim Manji-Jung\char95{}, Nov 19 2026\par
\%C A294422 The conjecture is true: the forcing term g(n) = a(n) - a(n-1) - a(n-2) is \char62{}= 0, so a grows exponentially and b(n) \char60{}= 2n+1 for large n; then e(n) = a(n) - phi*a(n-1) satisfies e(n) = -e(n-1)/phi + g(n) and is polynomially bounded since g(n) is, so a(n)/a(n-1) - phi = e(n)/a(n-1) -\char62{} 0. - \char95{}Joaquim Manji-Jung\char95{}, Nov 19 2026\par
\%C A294423 The conjecture is true: the forcing term g(n) = a(n) - a(n-1) - a(n-2) is \char62{}= 0, so a grows exponentially and b(n) \char60{}= 2n+1 for large n; then e(n) = a(n) - phi*a(n-1) satisfies e(n) = -e(n-1)/phi + g(n) and is polynomially bounded since g(n) is, so a(n)/a(n-1) - phi = e(n)/a(n-1) -\char62{} 0. - \char95{}Joaquim Manji-Jung\char95{}, Nov 19 2026\par
\%C A294424 The conjecture is true: the forcing term g(n) = a(n) - a(n-1) - a(n-2) is \char62{}= 0, so a grows exponentially and b(n) \char60{}= 2n+1 for large n; then e(n) = a(n) - phi*a(n-1) satisfies e(n) = -e(n-1)/phi + g(n) and is polynomially bounded since g(n) is, so a(n)/a(n-1) - phi = e(n)/a(n-1) -\char62{} 0. - \char95{}Joaquim Manji-Jung\char95{}, Nov 19 2026\par
\%C A294425 The conjecture is true: the forcing term g(n) = a(n) - a(n-1) - a(n-2) is \char62{}= 0, so a grows exponentially and b(n) \char60{}= 2n+1 for large n; then e(n) = a(n) - phi*a(n-1) satisfies e(n) = -e(n-1)/phi + g(n) and is polynomially bounded since g(n) is, so a(n)/a(n-1) - phi = e(n)/a(n-1) -\char62{} 0. - \char95{}Joaquim Manji-Jung\char95{}, Nov 19 2026\par
\%C A294426 The conjecture is true: the forcing term g(n) = a(n) - a(n-1) - a(n-2) is \char62{}= 0, so a grows exponentially and b(n) \char60{}= 2n+1 for large n; then e(n) = a(n) - phi*a(n-1) satisfies e(n) = -e(n-1)/phi + g(n) and is polynomially bounded since g(n) is, so a(n)/a(n-1) - phi = e(n)/a(n-1) -\char62{} 0. - \char95{}Joaquim Manji-Jung\char95{}, Nov 19 2026\par
\%C A294532 The conjecture is true: the forcing term g(n) = a(n) - a(n-1) - a(n-2) is \char62{}= 0, so a grows exponentially and b(n) \char60{}= 2n+1 for large n; then e(n) = a(n) - phi*a(n-1) satisfies e(n) = -e(n-1)/phi + g(n) and is polynomially bounded since g(n) is, so a(n)/a(n-1) - phi = e(n)/a(n-1) -\char62{} 0. - \char95{}Joaquim Manji-Jung\char95{}, Nov 19 2026\par
\%C A294533 The conjecture is true: the forcing term g(n) = a(n) - a(n-1) - a(n-2) is \char62{}= 0, so a grows exponentially and b(n) \char60{}= 2n+1 for large n; then e(n) = a(n) - phi*a(n-1) satisfies e(n) = -e(n-1)/phi + g(n) and is polynomially bounded since g(n) is, so a(n)/a(n-1) - phi = e(n)/a(n-1) -\char62{} 0. - \char95{}Joaquim Manji-Jung\char95{}, Nov 19 2026\par
\%C A294534 The conjecture is true: the forcing term g(n) = a(n) - a(n-1) - a(n-2) is \char62{}= 0, so a grows exponentially and b(n) \char60{}= 2n+1 for large n; then e(n) = a(n) - phi*a(n-1) satisfies e(n) = -e(n-1)/phi + g(n) and is polynomially bounded since g(n) is, so a(n)/a(n-1) - phi = e(n)/a(n-1) -\char62{} 0. - \char95{}Joaquim Manji-Jung\char95{}, Nov 19 2026\par
\%C A294535 The conjecture is true: the forcing term g(n) = a(n) - a(n-1) - a(n-2) is \char62{}= 0, so a grows exponentially and b(n) \char60{}= 2n+1 for large n; then e(n) = a(n) - phi*a(n-1) satisfies e(n) = -e(n-1)/phi + g(n) and is polynomially bounded since g(n) is, so a(n)/a(n-1) - phi = e(n)/a(n-1) -\char62{} 0. - \char95{}Joaquim Manji-Jung\char95{}, Nov 19 2026\par
\%C A294536 The conjecture is true: the forcing term g(n) = a(n) - a(n-1) - a(n-2) is \char62{}= 0, so a grows exponentially and b(n) \char60{}= 2n+1 for large n; then e(n) = a(n) - phi*a(n-1) satisfies e(n) = -e(n-1)/phi + g(n) and is polynomially bounded since g(n) is, so a(n)/a(n-1) - phi = e(n)/a(n-1) -\char62{} 0. - \char95{}Joaquim Manji-Jung\char95{}, Nov 19 2026\par
\%C A294537 The conjecture is true: the forcing term g(n) = a(n) - a(n-1) - a(n-2) is \char62{}= 0, so a grows exponentially and b(n) \char60{}= 2n+1 for large n; then e(n) = a(n) - phi*a(n-1) satisfies e(n) = -e(n-1)/phi + g(n) and is polynomially bounded since g(n) is, so a(n)/a(n-1) - phi = e(n)/a(n-1) -\char62{} 0. - \char95{}Joaquim Manji-Jung\char95{}, Nov 19 2026\par
\%C A294538 The conjecture is true: the forcing term g(n) = a(n) - a(n-1) - a(n-2) is \char62{}= 0, so a grows exponentially and b(n) \char60{}= 2n+1 for large n; then e(n) = a(n) - phi*a(n-1) satisfies e(n) = -e(n-1)/phi + g(n) and is polynomially bounded since g(n) is, so a(n)/a(n-1) - phi = e(n)/a(n-1) -\char62{} 0. - \char95{}Joaquim Manji-Jung\char95{}, Nov 19 2026\par
\%C A294539 The conjecture is true: the forcing term g(n) = a(n) - a(n-1) - a(n-2) is \char62{}= 0, so a grows exponentially and b(n) \char60{}= 2n+1 for large n; then e(n) = a(n) - phi*a(n-1) satisfies e(n) = -e(n-1)/phi + g(n) and is polynomially bounded since g(n) is, so a(n)/a(n-1) - phi = e(n)/a(n-1) -\char62{} 0. - \char95{}Joaquim Manji-Jung\char95{}, Nov 19 2026\par
\%C A294540 The conjecture is true: the forcing term g(n) = a(n) - a(n-1) - a(n-2) is \char62{}= 0, so a grows exponentially and b(n) \char60{}= 2n+1 for large n; then e(n) = a(n) - phi*a(n-1) satisfies e(n) = -e(n-1)/phi + g(n) and is polynomially bounded since g(n) is, so a(n)/a(n-1) - phi = e(n)/a(n-1) -\char62{} 0. - \char95{}Joaquim Manji-Jung\char95{}, Nov 19 2026\par
\%C A294541 The conjecture is true: the forcing term g(n) = a(n) - a(n-1) - a(n-2) is \char62{}= 0, so a grows exponentially and b(n) \char60{}= 2n+1 for large n; then e(n) = a(n) - phi*a(n-1) satisfies e(n) = -e(n-1)/phi + g(n) and is polynomially bounded since g(n) is, so a(n)/a(n-1) - phi = e(n)/a(n-1) -\char62{} 0. - \char95{}Joaquim Manji-Jung\char95{}, Nov 19 2026\par
\%C A294542 The conjecture is true: the forcing term g(n) = a(n) - a(n-1) - a(n-2) is \char62{}= 0, so a grows exponentially and b(n) \char60{}= 2n+1 for large n; then e(n) = a(n) - phi*a(n-1) satisfies e(n) = -e(n-1)/phi + g(n) and is polynomially bounded since g(n) is, so a(n)/a(n-1) - phi = e(n)/a(n-1) -\char62{} 0. - \char95{}Joaquim Manji-Jung\char95{}, Nov 19 2026\par
\%C A294543 The conjecture is true: the forcing term g(n) = a(n) - a(n-1) - a(n-2) is \char62{}= 0, so a grows exponentially and b(n) \char60{}= 2n+1 for large n; then e(n) = a(n) - phi*a(n-1) satisfies e(n) = -e(n-1)/phi + g(n) and is polynomially bounded since g(n) is, so a(n)/a(n-1) - phi = e(n)/a(n-1) -\char62{} 0. - \char95{}Joaquim Manji-Jung\char95{}, Nov 19 2026\par
\%C A294544 The conjecture is true: the forcing term g(n) = a(n) - a(n-1) - a(n-2) is \char62{}= 0, so a grows exponentially and b(n) \char60{}= 2n+1 for large n; then e(n) = a(n) - phi*a(n-1) satisfies e(n) = -e(n-1)/phi + g(n) and is polynomially bounded since g(n) is, so a(n)/a(n-1) - phi = e(n)/a(n-1) -\char62{} 0. - \char95{}Joaquim Manji-Jung\char95{}, Nov 19 2026\par
\%C A294545 The conjecture is true: the forcing term g(n) = a(n) - a(n-1) - a(n-2) is \char62{}= 0, so a grows exponentially and b(n) \char60{}= 2n+1 for large n; then e(n) = a(n) - phi*a(n-1) satisfies e(n) = -e(n-1)/phi + g(n) and is polynomially bounded since g(n) is, so a(n)/a(n-1) - phi = e(n)/a(n-1) -\char62{} 0. - \char95{}Joaquim Manji-Jung\char95{}, Nov 19 2026\par
\%C A294546 The conjecture is true: the forcing term g(n) = a(n) - a(n-1) - a(n-2) is \char62{}= 0, so a grows exponentially and b(n) \char60{}= 2n+1 for large n; then e(n) = a(n) - phi*a(n-1) satisfies e(n) = -e(n-1)/phi + g(n) and is polynomially bounded since g(n) is, so a(n)/a(n-1) - phi = e(n)/a(n-1) -\char62{} 0. - \char95{}Joaquim Manji-Jung\char95{}, Nov 19 2026\par
\%C A294547 The conjecture is true: the forcing term g(n) = a(n) - a(n-1) - a(n-2) is \char62{}= 0, so a grows exponentially and b(n) \char60{}= 2n+1 for large n; then e(n) = a(n) - phi*a(n-1) satisfies e(n) = -e(n-1)/phi + g(n) and is polynomially bounded since g(n) is, so a(n)/a(n-1) - phi = e(n)/a(n-1) -\char62{} 0. - \char95{}Joaquim Manji-Jung\char95{}, Nov 19 2026\par
\%C A294548 The conjecture is true: the forcing term g(n) = a(n) - a(n-1) - a(n-2) is \char62{}= 0, so a grows exponentially and b(n) \char60{}= 2n+1 for large n; then e(n) = a(n) - phi*a(n-1) satisfies e(n) = -e(n-1)/phi + g(n) and is polynomially bounded since g(n) is, so a(n)/a(n-1) - phi = e(n)/a(n-1) -\char62{} 0. - \char95{}Joaquim Manji-Jung\char95{}, Nov 19 2026\par
\%C A294549 The conjecture is true: the forcing term g(n) = a(n) - a(n-1) - a(n-2) is \char62{}= 0, so a grows exponentially and b(n) \char60{}= 2n+1 for large n; then e(n) = a(n) - phi*a(n-1) satisfies e(n) = -e(n-1)/phi + g(n) and is polynomially bounded since g(n) is, so a(n)/a(n-1) - phi = e(n)/a(n-1) -\char62{} 0. - \char95{}Joaquim Manji-Jung\char95{}, Nov 19 2026\par
\%C A294550 The conjecture is true: the forcing term g(n) = a(n) - a(n-1) - a(n-2) is \char62{}= 0, so a grows exponentially and b(n) \char60{}= 2n+1 for large n; then e(n) = a(n) - phi*a(n-1) satisfies e(n) = -e(n-1)/phi + g(n) and is polynomially bounded since g(n) is, so a(n)/a(n-1) - phi = e(n)/a(n-1) -\char62{} 0. - \char95{}Joaquim Manji-Jung\char95{}, Nov 19 2026\par
\%C A294551 The conjecture is true: the forcing term g(n) = a(n) - a(n-1) - a(n-2) is \char62{}= 0, so a grows exponentially and b(n) \char60{}= 2n+1 for large n; then e(n) = a(n) - phi*a(n-1) satisfies e(n) = -e(n-1)/phi + g(n) and is polynomially bounded since g(n) is, so a(n)/a(n-1) - phi = e(n)/a(n-1) -\char62{} 0. - \char95{}Joaquim Manji-Jung\char95{}, Nov 19 2026\par
\%C A294552 Proof of the limit: the forcing term g(n) = a(n) - a(n-1) - a(n-2) is \char62{}= 0, so a grows exponentially and b(n) \char60{}= 2n+1 for large n; then e(n) = a(n) - phi*a(n-1) satisfies e(n) = -e(n-1)/phi + g(n) and is polynomially bounded since g(n) is, so a(n)/a(n-1) - phi = e(n)/a(n-1) -\char62{} 0. - \char95{}Joaquim Manji-Jung\char95{}, Nov 19 2026\par
\%C A294553 The conjecture is true: the forcing term g(n) = a(n) - a(n-1) - a(n-2) is \char62{}= 0, so a grows exponentially and b(n) \char60{}= 2n+1 for large n; then e(n) = a(n) - phi*a(n-1) satisfies e(n) = -e(n-1)/phi + g(n) and is polynomially bounded since g(n) is, so a(n)/a(n-1) - phi = e(n)/a(n-1) -\char62{} 0. - \char95{}Joaquim Manji-Jung\char95{}, Nov 19 2026\par
\%C A294554 The conjecture is true: the forcing term g(n) = a(n) - a(n-1) - a(n-2) is \char62{}= 0, so a grows exponentially and b(n) \char60{}= 2n+1 for large n; then e(n) = a(n) - phi*a(n-1) satisfies e(n) = -e(n-1)/phi + g(n) and is polynomially bounded since g(n) is, so a(n)/a(n-1) - phi = e(n)/a(n-1) -\char62{} 0. - \char95{}Joaquim Manji-Jung\char95{}, Nov 19 2026\par
\%C A294555 The conjecture is true: the forcing term g(n) = a(n) - a(n-1) - a(n-2) is \char62{}= 0, so a grows exponentially and b(n) \char60{}= 2n+1 for large n; then e(n) = a(n) - phi*a(n-1) satisfies e(n) = -e(n-1)/phi + g(n) and is polynomially bounded since g(n) is, so a(n)/a(n-1) - phi = e(n)/a(n-1) -\char62{} 0. - \char95{}Joaquim Manji-Jung\char95{}, Nov 19 2026\par
\%C A294556 The conjecture is true: the forcing term g(n) = a(n) - a(n-1) - a(n-2) is \char62{}= 0, so a grows exponentially and b(n) \char60{}= 2n+1 for large n; then e(n) = a(n) - phi*a(n-1) satisfies e(n) = -e(n-1)/phi + g(n) and is polynomially bounded since g(n) is, so a(n)/a(n-1) - phi = e(n)/a(n-1) -\char62{} 0. - \char95{}Joaquim Manji-Jung\char95{}, Nov 19 2026\par
\%C A294557 The conjecture is true: the forcing term g(n) = a(n) - a(n-1) - a(n-2) is \char62{}= 0, so a grows exponentially and b(n) \char60{}= 2n+1 for large n; then e(n) = a(n) - phi*a(n-1) satisfies e(n) = -e(n-1)/phi + g(n) and is polynomially bounded since g(n) is, so a(n)/a(n-1) - phi = e(n)/a(n-1) -\char62{} 0. - \char95{}Joaquim Manji-Jung\char95{}, Nov 19 2026\par
\%C A294558 The conjecture is true: the forcing term g(n) = a(n) - a(n-1) - a(n-2) is \char62{}= 0, so a grows exponentially and b(n) \char60{}= 2n+1 for large n; then e(n) = a(n) - phi*a(n-1) satisfies e(n) = -e(n-1)/phi + g(n) and is polynomially bounded since g(n) is, so a(n)/a(n-1) - phi = e(n)/a(n-1) -\char62{} 0. - \char95{}Joaquim Manji-Jung\char95{}, Nov 19 2026\par
\%C A294559 The conjecture is true: the forcing term g(n) = a(n) - a(n-1) - a(n-2) is \char62{}= 0, so a grows exponentially and b(n) \char60{}= 2n+1 for large n; then e(n) = a(n) - phi*a(n-1) satisfies e(n) = -e(n-1)/phi + g(n) and is polynomially bounded since g(n) is, so a(n)/a(n-1) - phi = e(n)/a(n-1) -\char62{} 0. - \char95{}Joaquim Manji-Jung\char95{}, Nov 19 2026\par
\%C A294560 The conjecture is true: the forcing term g(n) = a(n) - a(n-1) - a(n-2) is \char62{}= 0, so a grows exponentially and b(n) \char60{}= 2n+1 for large n; then e(n) = a(n) - phi*a(n-1) satisfies e(n) = -e(n-1)/phi + g(n) and is polynomially bounded since g(n) is, so a(n)/a(n-1) - phi = e(n)/a(n-1) -\char62{} 0. - \char95{}Joaquim Manji-Jung\char95{}, Nov 19 2026\par
\%C A294561 The conjecture is true: the forcing term g(n) = a(n) - a(n-1) - a(n-2) is \char62{}= 0, so a grows exponentially and b(n) \char60{}= 2n+1 for large n; then e(n) = a(n) - phi*a(n-1) satisfies e(n) = -e(n-1)/phi + g(n) and is polynomially bounded since g(n) is, so a(n)/a(n-1) - phi = e(n)/a(n-1) -\char62{} 0. - \char95{}Joaquim Manji-Jung\char95{}, Nov 19 2026\par
\%C A294562 The conjecture is true: the forcing term g(n) = a(n) - a(n-1) - a(n-2) is \char62{}= 0, so a grows exponentially and b(n) \char60{}= 2n+1 for large n; then e(n) = a(n) - phi*a(n-1) satisfies e(n) = -e(n-1)/phi + g(n) and is polynomially bounded since g(n) is, so a(n)/a(n-1) - phi = e(n)/a(n-1) -\char62{} 0. - \char95{}Joaquim Manji-Jung\char95{}, Nov 19 2026\par
\%C A294563 The conjecture is true: the forcing term g(n) = a(n) - a(n-1) - a(n-2) is \char62{}= 0, so a grows exponentially and b(n) \char60{}= 2n+1 for large n; then e(n) = a(n) - phi*a(n-1) satisfies e(n) = -e(n-1)/phi + g(n) and is polynomially bounded since g(n) is, so a(n)/a(n-1) - phi = e(n)/a(n-1) -\char62{} 0. - \char95{}Joaquim Manji-Jung\char95{}, Nov 19 2026\par
\%C A294564 The conjecture is true: the forcing term g(n) = a(n) - a(n-1) - a(n-2) is \char62{}= 0, so a grows exponentially and b(n) \char60{}= 2n+1 for large n; then e(n) = a(n) - phi*a(n-1) satisfies e(n) = -e(n-1)/phi + g(n) and is polynomially bounded since g(n) is, so a(n)/a(n-1) - phi = e(n)/a(n-1) -\char62{} 0. - \char95{}Joaquim Manji-Jung\char95{}, Nov 19 2026\par
\%C A294565 The conjecture is true: the forcing term g(n) = a(n) - a(n-1) - a(n-2) is \char62{}= 0, so a grows exponentially and b(n) \char60{}= 2n+1 for large n; then e(n) = a(n) - phi*a(n-1) satisfies e(n) = -e(n-1)/phi + g(n) and is polynomially bounded since g(n) is, so a(n)/a(n-1) - phi = e(n)/a(n-1) -\char62{} 0. - \char95{}Joaquim Manji-Jung\char95{}, Nov 19 2026\par
\%C A295357 Proof of the limit: the forcing term g(n) = a(n) - a(n-1) - a(n-2) is \char62{}= 0, so a grows exponentially and b(n) \char60{}= 2n+1 for large n; then e(n) = a(n) - phi*a(n-1) satisfies e(n) = -e(n-1)/phi + g(n) and is polynomially bounded since g(n) is, so a(n)/a(n-1) - phi = e(n)/a(n-1) -\char62{} 0. - \char95{}Joaquim Manji-Jung\char95{}, Nov 19 2026\par
\%C A295754 Proof of the limit: b(n-4) \char62{}= 0 gives a(n) \char62{}= a(n-1) + a(n-3), so a grows exponentially and b(n) \char60{}= 2n+1 for large n; the characteristic polynomial is (x\char94{}2-x-1)(x\char94{}2+1), so e(n) = a(n) - phi*a(n-1) satisfies a recurrence whose roots -1/phi, i, -i are simple of modulus \char60{}= 1, forced by b(n-4) = O(n); hence e(n) = O(n\char94{}2) and a(n)/a(n-1) -\char62{} phi. - \char95{}Joaquim Manji-Jung\char95{}, Nov 19 2026\par
\%C A295862 Proof of the limit: the forcing term g(n) = a(n) - a(n-1) - a(n-2) is \char62{}= 0, so a grows exponentially and b(n) \char60{}= 2n+1 for large n; then e(n) = a(n) - phi*a(n-1) satisfies e(n) = -e(n-1)/phi + g(n) and is polynomially bounded since g(n) is, so a(n)/a(n-1) - phi = e(n)/a(n-1) -\char62{} 0. - \char95{}Joaquim Manji-Jung\char95{}, Nov 19 2026\par
\%C A295947 Proof of the limit: the forcing term g(n) = a(n) - a(n-1) - a(n-2) is \char62{}= 0, so a grows exponentially and b(n) \char60{}= 2n+1 for large n; then e(n) = a(n) - phi*a(n-1) satisfies e(n) = -e(n-1)/phi + g(n) and is polynomially bounded since g(n) is, so a(n)/a(n-1) - phi = e(n)/a(n-1) -\char62{} 0. - \char95{}Joaquim Manji-Jung\char95{}, Nov 19 2026\par
\%C A295948 Proof of the limit: the forcing term g(n) = a(n) - a(n-1) - a(n-2) is \char62{}= 0, so a grows exponentially and b(n) \char60{}= 2n+1 for large n; then e(n) = a(n) - phi*a(n-1) satisfies e(n) = -e(n-1)/phi + g(n) and is polynomially bounded since g(n) is, so a(n)/a(n-1) - phi = e(n)/a(n-1) -\char62{} 0. - \char95{}Joaquim Manji-Jung\char95{}, Nov 19 2026\par
\%C A295949 Proof of the limit: the forcing term g(n) = a(n) - a(n-1) - a(n-2) is \char62{}= 0, so a grows exponentially and b(n) \char60{}= 2n+1 for large n; then e(n) = a(n) - phi*a(n-1) satisfies e(n) = -e(n-1)/phi + g(n) and is polynomially bounded since g(n) is, so a(n)/a(n-1) - phi = e(n)/a(n-1) -\char62{} 0. - \char95{}Joaquim Manji-Jung\char95{}, Nov 19 2026\par
\%C A295950 Proof of the limit: the forcing term g(n) = a(n) - a(n-1) - a(n-2) is \char62{}= 0, so a grows exponentially and b(n) \char60{}= 2n+1 for large n; then e(n) = a(n) - phi*a(n-1) satisfies e(n) = -e(n-1)/phi + g(n) and is polynomially bounded since g(n) is, so a(n)/a(n-1) - phi = e(n)/a(n-1) -\char62{} 0. - \char95{}Joaquim Manji-Jung\char95{}, Nov 19 2026\par
\%C A295951 Proof of the limit: the forcing term g(n) = a(n) - a(n-1) - a(n-2) is \char62{}= 0, so a grows exponentially and b(n) \char60{}= 2n+1 for large n; then e(n) = a(n) - phi*a(n-1) satisfies e(n) = -e(n-1)/phi + g(n) and is polynomially bounded since g(n) is, so a(n)/a(n-1) - phi = e(n)/a(n-1) -\char62{} 0. - \char95{}Joaquim Manji-Jung\char95{}, Nov 19 2026\par
\%C A295952 Proof of the limit: the forcing term g(n) = a(n) - a(n-1) - a(n-2) is \char62{}= 0, so a grows exponentially and b(n) \char60{}= 2n+1 for large n; then e(n) = a(n) - phi*a(n-1) satisfies e(n) = -e(n-1)/phi + g(n) and is polynomially bounded since g(n) is, so a(n)/a(n-1) - phi = e(n)/a(n-1) -\char62{} 0. - \char95{}Joaquim Manji-Jung\char95{}, Nov 19 2026\par
\%C A295953 Proof of the limit: the forcing term g(n) = a(n) - a(n-1) - a(n-2) is \char62{}= 0, so a grows exponentially and b(n) \char60{}= 2n+1 for large n; then e(n) = a(n) - phi*a(n-1) satisfies e(n) = -e(n-1)/phi + g(n) and is polynomially bounded since g(n) is, so a(n)/a(n-1) - phi = e(n)/a(n-1) -\char62{} 0. - \char95{}Joaquim Manji-Jung\char95{}, Nov 19 2026\par
\%C A295960 Proof of the limit: the forcing term g(n) = a(n) - a(n-1) - a(n-2) is \char62{}= 0, so a grows exponentially and b(n) \char60{}= 2n+1 for large n; then e(n) = a(n) - phi*a(n-1) satisfies e(n) = -e(n-1)/phi + g(n) and is polynomially bounded since g(n) is, so a(n)/a(n-1) - phi = e(n)/a(n-1) -\char62{} 0. - \char95{}Joaquim Manji-Jung\char95{}, Nov 19 2026\par
\%C A296245 Proof of the limit: the forcing term g(n) = a(n) - a(n-1) - a(n-2) is \char62{}= 0, so a grows exponentially and b(n) \char60{}= 2n+1 for large n; then e(n) = a(n) - phi*a(n-1) satisfies e(n) = -e(n-1)/phi + g(n) and is polynomially bounded since g(n) is, so a(n)/a(n-1) - phi = e(n)/a(n-1) -\char62{} 0. - \char95{}Joaquim Manji-Jung\char95{}, Nov 19 2026\par
\%C A296251 Proof of the limit: the forcing term g(n) = a(n) - a(n-1) - a(n-2) is \char62{}= 0, so a grows exponentially and b(n) \char60{}= 2n+1 for large n; then e(n) = a(n) - phi*a(n-1) satisfies e(n) = -e(n-1)/phi + g(n) and is polynomially bounded since g(n) is, so a(n)/a(n-1) - phi = e(n)/a(n-1) -\char62{} 0. - \char95{}Joaquim Manji-Jung\char95{}, Nov 19 2026\par
\%C A296257 Proof of the limit: the forcing term g(n) = a(n) - a(n-1) - a(n-2) is \char62{}= 0, so a grows exponentially and b(n) \char60{}= 2n+1 for large n; then e(n) = a(n) - phi*a(n-1) satisfies e(n) = -e(n-1)/phi + g(n) and is polynomially bounded since g(n) is, so a(n)/a(n-1) - phi = e(n)/a(n-1) -\char62{} 0. - \char95{}Joaquim Manji-Jung\char95{}, Nov 19 2026\par
\%C A296266 Proof of the limit: the forcing term g(n) = a(n) - a(n-1) - a(n-2) is \char62{}= 0, so a grows exponentially and b(n) \char60{}= 2n+1 for large n; then e(n) = a(n) - phi*a(n-1) satisfies e(n) = -e(n-1)/phi + g(n) and is polynomially bounded since g(n) is, so a(n)/a(n-1) - phi = e(n)/a(n-1) -\char62{} 0. - \char95{}Joaquim Manji-Jung\char95{}, Nov 19 2026\par
\%C A296272 Proof of the limit: the forcing term g(n) = a(n) - a(n-1) - a(n-2) is \char62{}= 0, so a grows exponentially and b(n) \char60{}= 2n+1 for large n; then e(n) = a(n) - phi*a(n-1) satisfies e(n) = -e(n-1)/phi + g(n) and is polynomially bounded since g(n) is, so a(n)/a(n-1) - phi = e(n)/a(n-1) -\char62{} 0. - \char95{}Joaquim Manji-Jung\char95{}, Nov 19 2026\par
\%C A296278 Proof of the limit: the forcing term g(n) = a(n) - a(n-1) - a(n-2) is \char62{}= 0, so a grows exponentially and b(n) \char60{}= 2n+1 for large n; then e(n) = a(n) - phi*a(n-1) satisfies e(n) = -e(n-1)/phi + g(n) and is polynomially bounded since g(n) is, so a(n)/a(n-1) - phi = e(n)/a(n-1) -\char62{} 0. - \char95{}Joaquim Manji-Jung\char95{}, Nov 19 2026\par
\%C A296284 Proof of the limit: the forcing term g(n) = a(n) - a(n-1) - a(n-2) is \char62{}= 0, so a grows exponentially and b(n) \char60{}= 2n+1 for large n; then e(n) = a(n) - phi*a(n-1) satisfies e(n) = -e(n-1)/phi + g(n) and is polynomially bounded since g(n) is, so a(n)/a(n-1) - phi = e(n)/a(n-1) -\char62{} 0. - \char95{}Joaquim Manji-Jung\char95{}, Nov 19 2026\par
\%C A296288 Proof of the limit: the forcing term g(n) = a(n) - a(n-1) - a(n-2) is \char62{}= 0, so a grows exponentially and b(n) \char60{}= 2n+1 for large n; then e(n) = a(n) - phi*a(n-1) satisfies e(n) = -e(n-1)/phi + g(n) and is polynomially bounded since g(n) is, so a(n)/a(n-1) - phi = e(n)/a(n-1) -\char62{} 0. - \char95{}Joaquim Manji-Jung\char95{}, Nov 19 2026\par
\%C A296292 Proof of the limit: the forcing term g(n) = a(n) - a(n-1) - a(n-2) is \char62{}= 0, so a grows exponentially and b(n) \char60{}= 2n+1 for large n; then e(n) = a(n) - phi*a(n-1) satisfies e(n) = -e(n-1)/phi + g(n) and is polynomially bounded since g(n) is, so a(n)/a(n-1) - phi = e(n)/a(n-1) -\char62{} 0. - \char95{}Joaquim Manji-Jung\char95{}, Nov 19 2026\par
\%C A296555 Proof of the limit: the forcing term g(n) = a(n) - a(n-1) - a(n-2) is \char62{}= 0, so a grows exponentially and b(n) \char60{}= 2n+1 for large n; then e(n) = a(n) - phi*a(n-1) satisfies e(n) = -e(n-1)/phi + g(n) and is polynomially bounded since g(n) is, so a(n)/a(n-1) - phi = e(n)/a(n-1) -\char62{} 0. - \char95{}Joaquim Manji-Jung\char95{}, Nov 19 2026\par
\%C A296776 Proof of the limit: the forcing term g(n) = a(n) - a(n-1) - a(n-2) is \char62{}= 0, so a grows exponentially and b(n) \char60{}= 2n+1 for large n; then e(n) = a(n) - phi*a(n-1) satisfies e(n) = -e(n-1)/phi + g(n) and is polynomially bounded since g(n) is, so a(n)/a(n-1) - phi = e(n)/a(n-1) -\char62{} 0. - \char95{}Joaquim Manji-Jung\char95{}, Nov 19 2026\par
\%C A296843 Proof of the limit: the forcing term g(n) = a(n) - a(n-1) - a(n-2) is \char62{}= 0, so a grows exponentially and b(n) \char60{}= 2n+1 for large n; then e(n) = a(n) - phi*a(n-1) satisfies e(n) = -e(n-1)/phi + g(n) and is polynomially bounded since g(n) is, so a(n)/a(n-1) - phi = e(n)/a(n-1) -\char62{} 0. - \char95{}Joaquim Manji-Jung\char95{}, Nov 19 2026\par
\%C A295070 The conjecture is true: a(n) - a(n-2) = b(n-1) + b(n-2) \char62{}= 2n-1 makes a grow quadratically, so b(m) \char60{}= 2m+1 (counting), hence 0 \char60{} a(n) - a(n-1) \char60{}= 4n and a(n)/a(n-1) = 1 + O(1/n). - \char95{}Joaquim Manji-Jung\char95{}, Nov 19 2026\par
}
\subsection*{\texttt{conjectures/dfinite/draft.txt}}
{\small\raggedright\ttfamily
\%C A026743 The conjectured recurrence is true: applying the corresponding differential operator to the closed-form g.f. gives a polynomial (exact computation in Q(x)[sqrt(Delta)]), and the first terms are checked directly. - \char95{}Joaquim Manji-Jung\char95{}, Nov 19 2026\par
\%C A027174 The conjectured recurrence is true: applying the corresponding differential operator to the closed-form g.f. gives a polynomial (exact computation in Q(x)[sqrt(Delta)]), and the first terms are checked directly. - \char95{}Joaquim Manji-Jung\char95{}, Nov 19 2026\par
\%C A027175 The conjectured recurrence is true: applying the corresponding differential operator to the closed-form g.f. gives a polynomial (exact computation in Q(x)[sqrt(Delta)]), and the first terms are checked directly. - \char95{}Joaquim Manji-Jung\char95{}, Nov 19 2026\par
\%C A103973 The conjectured recurrence is true: applying the corresponding differential operator to the closed-form g.f. gives a polynomial (exact computation in Q(x)[sqrt(Delta)]), and the first terms are checked directly. - \char95{}Joaquim Manji-Jung\char95{}, Nov 19 2026\par
\%C A190171 The conjectured recurrence is true: applying the corresponding differential operator to the closed-form g.f. gives a polynomial (exact computation in Q(x)[sqrt(Delta)]), and the first terms are checked directly. - \char95{}Joaquim Manji-Jung\char95{}, Nov 19 2026\par
\%C A385252 The conjectured recurrence holds for n \char62{}= 6 (exact computation with the closed-form g.f.), but not for n = 5. - \char95{}Joaquim Manji-Jung\char95{}, Nov 19 2026\par
}
\subsection*{\texttt{conjectures/fibtransform/draft.txt}}
{\small\raggedright\ttfamily
\%C A287676 Kimberling's conjecture is true: in Zeckendorf numeration the positions of the letters of this transform of the Fibonacci word are synchronized (the block of the i-th letter of A003849 starts at 1 + i + 2*(2i + 1 - floor((i+1)*phi))), r lies in Q(phi), and the inequalities are proved with the Walnut theorem prover. - \char95{}Joaquim Manji-Jung\char95{}, Nov 19 2026\par
\%C A287723 Kimberling's conjecture is true: in Zeckendorf numeration the positions of the letters of this transform of the Fibonacci word are synchronized (the block of the i-th letter of A003849 starts at 1 + i + 2*(2i + 1 - floor((i+1)*phi))), r lies in Q(phi), and the inequalities are proved with the Walnut theorem prover. - \char95{}Joaquim Manji-Jung\char95{}, Nov 19 2026\par
\%C A287724 Kimberling's conjecture is true: in Zeckendorf numeration the positions of the letters of this transform of the Fibonacci word are synchronized (the block of the i-th letter of A003849 starts at 1 + i + 2*(2i + 1 - floor((i+1)*phi))), r lies in Q(phi), and the inequalities are proved with the Walnut theorem prover. - \char95{}Joaquim Manji-Jung\char95{}, Nov 19 2026\par
\%C A287727 Kimberling's conjecture is true: in Zeckendorf numeration the positions of the letters of this transform of the Fibonacci word are synchronized (the block of the i-th letter of A003849 starts at 1 + i + 2*(2i + 1 - floor((i+1)*phi))), r lies in Q(phi), and the inequalities are proved with the Walnut theorem prover. - \char95{}Joaquim Manji-Jung\char95{}, Nov 19 2026\par
\%C A287676 Chunqing Liu's conjectured formula a(n) = 5n + 2 - 2*floor((n+1)*phi) is true (proved with Walnut, see above); it gives n*r - a(n) = sqrt(5) - 1 - 2*frac((n+1)*phi). - \char95{}Joaquim Manji-Jung\char95{}, Nov 19 2026\par
}
\subsection*{\texttt{conjectures/kimberling/draft.txt}}
{\small\raggedright\ttfamily
\%C A026363 Kimberling's conjecture is true; the exact infimum and supremum of n*r - a(n) are computed in the companion repository (exact computation in Q(sqrt(d)) via Dumont-Thomas numeration, with a verified certificate). - \char95{}Joaquim Manji-Jung\char95{}, Nov 19 2026\par
\%C A026364 Kimberling's conjecture is true; the exact infimum and supremum of n*r - a(n) are computed in the companion repository (exact computation in Q(sqrt(d)) via Dumont-Thomas numeration, with a verified certificate). - \char95{}Joaquim Manji-Jung\char95{}, Nov 19 2026\par
\%C A026368 Kimberling's conjecture is true; the exact infimum and supremum of n*r - a(n) are computed in the companion repository (exact computation in Q(sqrt(d)) via Dumont-Thomas numeration, with a verified certificate). - \char95{}Joaquim Manji-Jung\char95{}, Nov 19 2026\par
\%C A284753 Kimberling's conjecture is true; the exact infimum and supremum of n*r - a(n) are computed in the companion repository (exact computation in Q(sqrt(d)) via Dumont-Thomas numeration, with a verified certificate). - \char95{}Joaquim Manji-Jung\char95{}, Nov 19 2026\par
\%C A284894 Kimberling's conjecture is true; the exact infimum and supremum of n*r - a(n) are computed in the companion repository (exact computation in Q(sqrt(d)) via Dumont-Thomas numeration, with a verified certificate). - \char95{}Joaquim Manji-Jung\char95{}, Nov 19 2026\par
\%C A285669 Kimberling's conjecture is true; the exact infimum and supremum of n*r - a(n) are computed in the companion repository (exact computation in Q(sqrt(d)) via Dumont-Thomas numeration, with a verified certificate). - \char95{}Joaquim Manji-Jung\char95{}, Nov 19 2026\par
\%C A026367 Kimberling's conjecture is true; the exact infimum and supremum of n*r - a(n) are computed in the companion repository (exact computation in Q(sqrt(d)) via Dumont-Thomas numeration, with a verified certificate). - \char95{}Joaquim Manji-Jung\char95{}, Nov 19 2026\par
\%C A284654 Kimberling's conjecture is true; the exact infimum and supremum of n*r - a(n) are computed in the companion repository (exact computation in Q(sqrt(d)) via Dumont-Thomas numeration, with a verified certificate). - \char95{}Joaquim Manji-Jung\char95{}, Nov 19 2026\par
\%C A284675 Kimberling's conjecture is true; the exact infimum and supremum of n*r - a(n) are computed in the companion repository (exact computation in Q(sqrt(d)) via Dumont-Thomas numeration, with a verified certificate). - \char95{}Joaquim Manji-Jung\char95{}, Nov 19 2026\par
\%C A284676 Kimberling's conjecture is true; the exact infimum and supremum of n*r - a(n) are computed in the companion repository (exact computation in Q(sqrt(d)) via Dumont-Thomas numeration, with a verified certificate). - \char95{}Joaquim Manji-Jung\char95{}, Nov 19 2026\par
\%C A284678 Kimberling's conjecture is true; the exact infimum and supremum of n*r - a(n) are computed in the companion repository (exact computation in Q(sqrt(d)) via Dumont-Thomas numeration, with a verified certificate). - \char95{}Joaquim Manji-Jung\char95{}, Nov 19 2026\par
\%C A284679 Kimberling's conjecture is true; the exact infimum and supremum of n*r - a(n) are computed in the companion repository (exact computation in Q(sqrt(d)) via Dumont-Thomas numeration, with a verified certificate). - \char95{}Joaquim Manji-Jung\char95{}, Nov 19 2026\par
\%C A285032 Kimberling's conjecture is true; the exact infimum and supremum of n*r - a(n) are computed in the companion repository (exact computation in Q(sqrt(d)) via Dumont-Thomas numeration, with a verified certificate). - \char95{}Joaquim Manji-Jung\char95{}, Nov 19 2026\par
\%C A285035 Kimberling's conjecture is true; the exact infimum and supremum of n*r - a(n) are computed in the companion repository (exact computation in Q(sqrt(d)) via Dumont-Thomas numeration, with a verified certificate). - \char95{}Joaquim Manji-Jung\char95{}, Nov 19 2026\par
\%C A285036 Kimberling's conjecture is true; the exact infimum and supremum of n*r - a(n) are computed in the companion repository (exact computation in Q(sqrt(d)) via Dumont-Thomas numeration, with a verified certificate). - \char95{}Joaquim Manji-Jung\char95{}, Nov 19 2026\par
\%C A285082 Kimberling's conjecture is true; the exact infimum and supremum of n*r - a(n) are computed in the companion repository (exact computation in Q(sqrt(d)) via Dumont-Thomas numeration, with a verified certificate). - \char95{}Joaquim Manji-Jung\char95{}, Nov 19 2026\par
\%C A285084 Kimberling's conjecture is true; the exact infimum and supremum of n*r - a(n) are computed in the companion repository (exact computation in Q(sqrt(d)) via Dumont-Thomas numeration, with a verified certificate). - \char95{}Joaquim Manji-Jung\char95{}, Nov 19 2026\par
\%C A285275 Kimberling's conjecture is true; the exact infimum and supremum of n*r - a(n) are computed in the companion repository (exact computation in Q(sqrt(d)) via Dumont-Thomas numeration, with a verified certificate). - \char95{}Joaquim Manji-Jung\char95{}, Nov 19 2026\par
\%C A285278 Kimberling's conjecture is true; the exact infimum and supremum of n*r - a(n) are computed in the companion repository (exact computation in Q(sqrt(d)) via Dumont-Thomas numeration, with a verified certificate). - \char95{}Joaquim Manji-Jung\char95{}, Nov 19 2026\par
\%C A285419 Kimberling's conjecture is true; the exact infimum and supremum of n*r - a(n) are computed in the companion repository (exact computation in Q(sqrt(d)) via Dumont-Thomas numeration, with a verified certificate). - \char95{}Joaquim Manji-Jung\char95{}, Nov 19 2026\par
\%C A283964 Kimberling's conjecture is true; the exact infimum and supremum of n*r - a(n) are computed in the companion repository (exact computation in Q(sqrt(d)) via Dumont-Thomas numeration, with a verified certificate). - \char95{}Joaquim Manji-Jung\char95{}, Nov 19 2026\par
\%C A283967 Kimberling's conjecture is true; the exact infimum and supremum of n*r - a(n) are computed in the companion repository (exact computation in Q(sqrt(d)) via Dumont-Thomas numeration, with a verified certificate). - \char95{}Joaquim Manji-Jung\char95{}, Nov 19 2026\par
\%C A284370 Kimberling's conjecture is true; the exact infimum and supremum of n*r - a(n) are computed in the companion repository (exact computation in Q(sqrt(d)) via Dumont-Thomas numeration, with a verified certificate). - \char95{}Joaquim Manji-Jung\char95{}, Nov 19 2026\par
\%C A284506 Kimberling's conjecture is true; the exact infimum and supremum of n*r - a(n) are computed in the companion repository (exact computation in Q(sqrt(d)) via Dumont-Thomas numeration, with a verified certificate). - \char95{}Joaquim Manji-Jung\char95{}, Nov 19 2026\par
\%C A284657 Kimberling's conjecture is true; the exact infimum and supremum of n*r - a(n) are computed in the companion repository (exact computation in Q(sqrt(d)) via Dumont-Thomas numeration, with a verified certificate). - \char95{}Joaquim Manji-Jung\char95{}, Nov 19 2026\par
\%C A284752 Kimberling's conjecture is true; the exact infimum and supremum of n*r - a(n) are computed in the companion repository (exact computation in Q(sqrt(d)) via Dumont-Thomas numeration, with a verified certificate). - \char95{}Joaquim Manji-Jung\char95{}, Nov 19 2026\par
\%C A284773 Kimberling's conjecture is true; the exact infimum and supremum of n*r - a(n) are computed in the companion repository (exact computation in Q(sqrt(d)) via Dumont-Thomas numeration, with a verified certificate). - \char95{}Joaquim Manji-Jung\char95{}, Nov 19 2026\par
\%C A284774 Kimberling's conjecture is true; the exact infimum and supremum of n*r - a(n) are computed in the companion repository (exact computation in Q(sqrt(d)) via Dumont-Thomas numeration, with a verified certificate). - \char95{}Joaquim Manji-Jung\char95{}, Nov 19 2026\par
\%C A284852 Kimberling's conjecture is true; the exact infimum and supremum of n*r - a(n) are computed in the companion repository (exact computation in Q(sqrt(d)) via Dumont-Thomas numeration, with a verified certificate). - \char95{}Joaquim Manji-Jung\char95{}, Nov 19 2026\par
\%C A284853 Kimberling's conjecture is true; the exact infimum and supremum of n*r - a(n) are computed in the companion repository (exact computation in Q(sqrt(d)) via Dumont-Thomas numeration, with a verified certificate). - \char95{}Joaquim Manji-Jung\char95{}, Nov 19 2026\par
\%C A284902 Kimberling's conjecture is true; the exact infimum and supremum of n*r - a(n) are computed in the companion repository (exact computation in Q(sqrt(d)) via Dumont-Thomas numeration, with a verified certificate). - \char95{}Joaquim Manji-Jung\char95{}, Nov 19 2026\par
\%C A284903 Kimberling's conjecture is true; the exact infimum and supremum of n*r - a(n) are computed in the companion repository (exact computation in Q(sqrt(d)) via Dumont-Thomas numeration, with a verified certificate). - \char95{}Joaquim Manji-Jung\char95{}, Nov 19 2026\par
\%C A284930 Kimberling's conjecture is true; the exact infimum and supremum of n*r - a(n) are computed in the companion repository (exact computation in Q(sqrt(d)) via Dumont-Thomas numeration, with a verified certificate). - \char95{}Joaquim Manji-Jung\char95{}, Nov 19 2026\par
\%C A284945 Kimberling's conjecture is true; the exact infimum and supremum of n*r - a(n) are computed in the companion repository (exact computation in Q(sqrt(d)) via Dumont-Thomas numeration, with a verified certificate). - \char95{}Joaquim Manji-Jung\char95{}, Nov 19 2026\par
\%C A284946 Kimberling's conjecture is true; the exact infimum and supremum of n*r - a(n) are computed in the companion repository (exact computation in Q(sqrt(d)) via Dumont-Thomas numeration, with a verified certificate). - \char95{}Joaquim Manji-Jung\char95{}, Nov 19 2026\par
\%C A285033 Kimberling's conjecture is true; the exact infimum and supremum of n*r - a(n) are computed in the companion repository (exact computation in Q(sqrt(d)) via Dumont-Thomas numeration, with a verified certificate). - \char95{}Joaquim Manji-Jung\char95{}, Nov 19 2026\par
\%C A285077 Kimberling's conjecture is true; the exact infimum and supremum of n*r - a(n) are computed in the companion repository (exact computation in Q(sqrt(d)) via Dumont-Thomas numeration, with a verified certificate). - \char95{}Joaquim Manji-Jung\char95{}, Nov 19 2026\par
\%C A285081 Kimberling's conjecture is true; the exact infimum and supremum of n*r - a(n) are computed in the companion repository (exact computation in Q(sqrt(d)) via Dumont-Thomas numeration, with a verified certificate). - \char95{}Joaquim Manji-Jung\char95{}, Nov 19 2026\par
\%C A285085 Kimberling's conjecture is true; the exact infimum and supremum of n*r - a(n) are computed in the companion repository (exact computation in Q(sqrt(d)) via Dumont-Thomas numeration, with a verified certificate). - \char95{}Joaquim Manji-Jung\char95{}, Nov 19 2026\par
\%C A285206 Kimberling's conjecture is true; the exact infimum and supremum of n*r - a(n) are computed in the companion repository (exact computation in Q(sqrt(d)) via Dumont-Thomas numeration, with a verified certificate). - \char95{}Joaquim Manji-Jung\char95{}, Nov 19 2026\par
\%C A285207 Kimberling's conjecture is true; the exact infimum and supremum of n*r - a(n) are computed in the companion repository (exact computation in Q(sqrt(d)) via Dumont-Thomas numeration, with a verified certificate). - \char95{}Joaquim Manji-Jung\char95{}, Nov 19 2026\par
\%C A285209 Kimberling's conjecture is true; the exact infimum and supremum of n*r - a(n) are computed in the companion repository (exact computation in Q(sqrt(d)) via Dumont-Thomas numeration, with a verified certificate). - \char95{}Joaquim Manji-Jung\char95{}, Nov 19 2026\par
\%C A285210 Kimberling's conjecture is true; the exact infimum and supremum of n*r - a(n) are computed in the companion repository (exact computation in Q(sqrt(d)) via Dumont-Thomas numeration, with a verified certificate). - \char95{}Joaquim Manji-Jung\char95{}, Nov 19 2026\par
\%C A285302 Kimberling's conjecture is true; the exact infimum and supremum of n*r - a(n) are computed in the companion repository (exact computation in Q(sqrt(d)) via Dumont-Thomas numeration, with a verified certificate). - \char95{}Joaquim Manji-Jung\char95{}, Nov 19 2026\par
\%C A285342 Kimberling's conjecture is true; the exact infimum and supremum of n*r - a(n) are computed in the companion repository (exact computation in Q(sqrt(d)) via Dumont-Thomas numeration, with a verified certificate). - \char95{}Joaquim Manji-Jung\char95{}, Nov 19 2026\par
\%C A285359 Kimberling's conjecture is true; the exact infimum and supremum of n*r - a(n) are computed in the companion repository (exact computation in Q(sqrt(d)) via Dumont-Thomas numeration, with a verified certificate). - \char95{}Joaquim Manji-Jung\char95{}, Nov 19 2026\par
\%C A285420 Kimberling's conjecture is true; the exact infimum and supremum of n*r - a(n) are computed in the companion repository (exact computation in Q(sqrt(d)) via Dumont-Thomas numeration, with a verified certificate). - \char95{}Joaquim Manji-Jung\char95{}, Nov 19 2026\par
\%C A285422 Kimberling's conjecture is true; the exact infimum and supremum of n*r - a(n) are computed in the companion repository (exact computation in Q(sqrt(d)) via Dumont-Thomas numeration, with a verified certificate). - \char95{}Joaquim Manji-Jung\char95{}, Nov 19 2026\par
\%C A285423 Kimberling's conjecture is true; the exact infimum and supremum of n*r - a(n) are computed in the companion repository (exact computation in Q(sqrt(d)) via Dumont-Thomas numeration, with a verified certificate). - \char95{}Joaquim Manji-Jung\char95{}, Nov 19 2026\par
\%C A285670 Kimberling's conjecture is true; the exact infimum and supremum of n*r - a(n) are computed in the companion repository (exact computation in Q(sqrt(d)) via Dumont-Thomas numeration, with a verified certificate). - \char95{}Joaquim Manji-Jung\char95{}, Nov 19 2026\par
\%C A285134 Kimberling's conjecture is false; see the counterexample in the companion repository (conjectures/kimberling/counter.out). - \char95{}Joaquim Manji-Jung\char95{}, Nov 19 2026\par
\%C A285135 Kimberling's conjecture is false; see the counterexample in the companion repository (conjectures/kimberling/counter.out). - \char95{}Joaquim Manji-Jung\char95{}, Nov 19 2026\par
\%C A285137 Kimberling's conjecture is false; see the counterexample in the companion repository (conjectures/kimberling/counter.out). - \char95{}Joaquim Manji-Jung\char95{}, Nov 19 2026\par
\%C A285138 Kimberling's conjecture is false; see the counterexample in the companion repository (conjectures/kimberling/counter.out). - \char95{}Joaquim Manji-Jung\char95{}, Nov 19 2026\par
\%C A285140 Kimberling's conjecture is false; see the counterexample in the companion repository (conjectures/kimberling/counter.out). - \char95{}Joaquim Manji-Jung\char95{}, Nov 19 2026\par
\%C A285141 Kimberling's conjecture is false; see the counterexample in the companion repository (conjectures/kimberling/counter.out). - \char95{}Joaquim Manji-Jung\char95{}, Nov 19 2026\par
\%C A285143 Kimberling's conjecture is false; see the counterexample in the companion repository (conjectures/kimberling/counter.out). - \char95{}Joaquim Manji-Jung\char95{}, Nov 19 2026\par
\%C A285144 Kimberling's conjecture is false; see the counterexample in the companion repository (conjectures/kimberling/counter.out). - \char95{}Joaquim Manji-Jung\char95{}, Nov 19 2026\par
\%C A285279 Kimberling's conjecture is false; see the counterexample in the companion repository (conjectures/kimberling/counter.out). - \char95{}Joaquim Manji-Jung\char95{}, Nov 19 2026\par
\%C A285129 Kimberling's conjecture is false; see the counterexample in the companion repository (conjectures/kimberling/counter.out). - \char95{}Joaquim Manji-Jung\char95{}, Nov 19 2026\par
\%C A285130 Kimberling's conjecture is false; see the counterexample in the companion repository (conjectures/kimberling/counter.out). - \char95{}Joaquim Manji-Jung\char95{}, Nov 19 2026\par
\%C A285078 The conjecture holds for n \char62{}= 2 with optimal bounds, but not for n = 1, where n*r - a(n) = sqrt(2) equals the upper bound. - \char95{}Joaquim Manji-Jung\char95{}, Nov 19 2026\par
\%C A284015 Kimberling's conjecture is true: the exact infimum and supremum of n*s - a(n) are (-7+sqrt(17))/4 and (1+sqrt(17))/4 (companion repository). - \char95{}Joaquim Manji-Jung\char95{}, Nov 19 2026\par
\%C A284469 The conjecture is true: the underlying morphism is primitive, so the letter has a frequency given by the Perron eigenvector, and its inverse is the stated limit (see the companion repository). - \char95{}Joaquim Manji-Jung\char95{}, Nov 19 2026\par
\%C A284470 The conjecture is true: the underlying morphism is primitive, so the letter has a frequency given by the Perron eigenvector, and its inverse is the stated limit (see the companion repository). - \char95{}Joaquim Manji-Jung\char95{}, Nov 19 2026\par
\%C A284472 The conjecture is true: the underlying morphism is primitive, so the letter has a frequency given by the Perron eigenvector, and its inverse is the stated limit (see the companion repository). - \char95{}Joaquim Manji-Jung\char95{}, Nov 19 2026\par
\%C A284473 The conjecture is true: the underlying morphism is primitive, so the letter has a frequency given by the Perron eigenvector, and its inverse is the stated limit (see the companion repository). - \char95{}Joaquim Manji-Jung\char95{}, Nov 19 2026\par
\%C A284509 The conjecture is true: the underlying morphism is primitive, so the letter has a frequency given by the Perron eigenvector, and its inverse is the stated limit (see the companion repository). - \char95{}Joaquim Manji-Jung\char95{}, Nov 19 2026\par
\%C A284510 The conjecture is true: the underlying morphism is primitive, so the letter has a frequency given by the Perron eigenvector, and its inverse is the stated limit (see the companion repository). - \char95{}Joaquim Manji-Jung\char95{}, Nov 19 2026\par
\%C A284513 The conjecture is true: the underlying morphism is primitive, so the letter has a frequency given by the Perron eigenvector, and its inverse is the stated limit (see the companion repository). - \char95{}Joaquim Manji-Jung\char95{}, Nov 19 2026\par
\%C A284514 The conjecture is true: the underlying morphism is primitive, so the letter has a frequency given by the Perron eigenvector, and its inverse is the stated limit (see the companion repository). - \char95{}Joaquim Manji-Jung\char95{}, Nov 19 2026\par
\%C A284525 The conjecture is true: the underlying morphism is primitive, so the letter has a frequency given by the Perron eigenvector, and its inverse is the stated limit (see the companion repository). - \char95{}Joaquim Manji-Jung\char95{}, Nov 19 2026\par
\%C A284526 The conjecture is true: the underlying morphism is primitive, so the letter has a frequency given by the Perron eigenvector, and its inverse is the stated limit (see the companion repository). - \char95{}Joaquim Manji-Jung\char95{}, Nov 19 2026\par
\%C A284528 The conjecture is true: the underlying morphism is primitive, so the letter has a frequency given by the Perron eigenvector, and its inverse is the stated limit (see the companion repository). - \char95{}Joaquim Manji-Jung\char95{}, Nov 19 2026\par
\%C A284529 The conjecture is true: the underlying morphism is primitive, so the letter has a frequency given by the Perron eigenvector, and its inverse is the stated limit (see the companion repository). - \char95{}Joaquim Manji-Jung\char95{}, Nov 19 2026\par
\%C A284534 The conjecture is true: the underlying morphism is primitive, so the letter has a frequency given by the Perron eigenvector, and its inverse is the stated limit (see the companion repository). - \char95{}Joaquim Manji-Jung\char95{}, Nov 19 2026\par
\%C A284535 The conjecture is true: the underlying morphism is primitive, so the letter has a frequency given by the Perron eigenvector, and its inverse is the stated limit (see the companion repository). - \char95{}Joaquim Manji-Jung\char95{}, Nov 19 2026\par
\%C A284589 The conjecture is true: the underlying morphism is primitive, so the letter has a frequency given by the Perron eigenvector, and its inverse is the stated limit (see the companion repository). - \char95{}Joaquim Manji-Jung\char95{}, Nov 19 2026\par
\%C A284590 The conjecture is true: the underlying morphism is primitive, so the letter has a frequency given by the Perron eigenvector, and its inverse is the stated limit (see the companion repository). - \char95{}Joaquim Manji-Jung\char95{}, Nov 19 2026\par
\%C A284681 The conjecture is true: the underlying morphism is primitive, so the letter has a frequency given by the Perron eigenvector, and its inverse is the stated limit (see the companion repository). - \char95{}Joaquim Manji-Jung\char95{}, Nov 19 2026\par
\%C A284682 The conjecture is true: the underlying morphism is primitive, so the letter has a frequency given by the Perron eigenvector, and its inverse is the stated limit (see the companion repository). - \char95{}Joaquim Manji-Jung\char95{}, Nov 19 2026\par
\%C A284955 The conjecture is true: the underlying morphism is primitive, so the letter has a frequency given by the Perron eigenvector, and its inverse is the stated limit (see the companion repository). - \char95{}Joaquim Manji-Jung\char95{}, Nov 19 2026\par
\%C A284956 The conjecture is true: the underlying morphism is primitive, so the letter has a frequency given by the Perron eigenvector, and its inverse is the stated limit (see the companion repository). - \char95{}Joaquim Manji-Jung\char95{}, Nov 19 2026\par
\%C A284958 The conjecture is true: the underlying morphism is primitive, so the letter has a frequency given by the Perron eigenvector, and its inverse is the stated limit (see the companion repository). - \char95{}Joaquim Manji-Jung\char95{}, Nov 19 2026\par
\%C A284959 The conjecture is true: the underlying morphism is primitive, so the letter has a frequency given by the Perron eigenvector, and its inverse is the stated limit (see the companion repository). - \char95{}Joaquim Manji-Jung\char95{}, Nov 19 2026\par
\%C A285160 The conjecture is true: the underlying morphism is primitive, so the letter has a frequency given by the Perron eigenvector, and its inverse is the stated limit (see the companion repository). - \char95{}Joaquim Manji-Jung\char95{}, Nov 19 2026\par
\%C A285161 The conjecture is true: the underlying morphism is primitive, so the letter has a frequency given by the Perron eigenvector, and its inverse is the stated limit (see the companion repository). - \char95{}Joaquim Manji-Jung\char95{}, Nov 19 2026\par
\%C A285502 The conjecture is true: the underlying morphism is primitive, so the letter has a frequency given by the Perron eigenvector, and its inverse is the stated limit (see the companion repository). - \char95{}Joaquim Manji-Jung\char95{}, Nov 19 2026\par
\%C A285505 The conjecture is true: the underlying morphism is primitive, so the letter has a frequency given by the Perron eigenvector, and its inverse is the stated limit (see the companion repository). - \char95{}Joaquim Manji-Jung\char95{}, Nov 19 2026\par
\%C A285506 The conjecture is true: the underlying morphism is primitive, so the letter has a frequency given by the Perron eigenvector, and its inverse is the stated limit (see the companion repository). - \char95{}Joaquim Manji-Jung\char95{}, Nov 19 2026\par
\%C A285516 The conjecture is true: the underlying morphism is primitive, so the letter has a frequency given by the Perron eigenvector, and its inverse is the stated limit (see the companion repository). - \char95{}Joaquim Manji-Jung\char95{}, Nov 19 2026\par
\%C A285531 The conjecture is true: the underlying morphism is primitive, so the letter has a frequency given by the Perron eigenvector, and its inverse is the stated limit (see the companion repository). - \char95{}Joaquim Manji-Jung\char95{}, Nov 19 2026\par
\%C A285532 The conjecture is true: the underlying morphism is primitive, so the letter has a frequency given by the Perron eigenvector, and its inverse is the stated limit (see the companion repository). - \char95{}Joaquim Manji-Jung\char95{}, Nov 19 2026\par
\%C A285534 The conjecture is true: the underlying morphism is primitive, so the letter has a frequency given by the Perron eigenvector, and its inverse is the stated limit (see the companion repository). - \char95{}Joaquim Manji-Jung\char95{}, Nov 19 2026\par
\%C A285535 The conjecture is true: the underlying morphism is primitive, so the letter has a frequency given by the Perron eigenvector, and its inverse is the stated limit (see the companion repository). - \char95{}Joaquim Manji-Jung\char95{}, Nov 19 2026\par
\%C A285566 The conjecture is true: the underlying morphism is primitive, so the letter has a frequency given by the Perron eigenvector, and its inverse is the stated limit (see the companion repository). - \char95{}Joaquim Manji-Jung\char95{}, Nov 19 2026\par
\%C A285569 The conjecture is true: the underlying morphism is primitive, so the letter has a frequency given by the Perron eigenvector, and its inverse is the stated limit (see the companion repository). - \char95{}Joaquim Manji-Jung\char95{}, Nov 19 2026\par
\%C A285597 The conjecture is true: the underlying morphism is primitive, so the letter has a frequency given by the Perron eigenvector, and its inverse is the stated limit (see the companion repository). - \char95{}Joaquim Manji-Jung\char95{}, Nov 19 2026\par
\%C A285598 The conjecture is true: the underlying morphism is primitive, so the letter has a frequency given by the Perron eigenvector, and its inverse is the stated limit (see the companion repository). - \char95{}Joaquim Manji-Jung\char95{}, Nov 19 2026\par
\%C A285601 The conjecture is true: the underlying morphism is primitive, so the letter has a frequency given by the Perron eigenvector, and its inverse is the stated limit (see the companion repository). - \char95{}Joaquim Manji-Jung\char95{}, Nov 19 2026\par
\%C A285619 The conjecture is true: the underlying morphism is primitive, so the letter has a frequency given by the Perron eigenvector, and its inverse is the stated limit (see the companion repository). - \char95{}Joaquim Manji-Jung\char95{}, Nov 19 2026\par
\%C A285658 The conjecture is true: the underlying morphism is primitive, so the letter has a frequency given by the Perron eigenvector, and its inverse is the stated limit (see the companion repository). - \char95{}Joaquim Manji-Jung\char95{}, Nov 19 2026\par
\%C A285659 The conjecture is true: the underlying morphism is primitive, so the letter has a frequency given by the Perron eigenvector, and its inverse is the stated limit (see the companion repository). - \char95{}Joaquim Manji-Jung\char95{}, Nov 19 2026\par
\%C A284488 The conjectured limit is incorrect: a(n)/n tends to 1/f, where f is the frequency of the letter given by the Perron eigenvector of the underlying primitive morphism, namely sqrt(3)/2 + 3/2 $\approx$ 2.366025. - \char95{}Joaquim Manji-Jung\char95{}, Nov 19 2026\par
\%C A284489 The conjectured limit is incorrect: a(n)/n tends to 1/f, where f is the frequency of the letter given by the Perron eigenvector of the underlying primitive morphism, namely sqrt(3) $\approx$ 1.732051. - \char95{}Joaquim Manji-Jung\char95{}, Nov 19 2026\par
\%C A284491 The conjectured limit is incorrect: a(n)/n tends to 1/f, where f is the frequency of the letter given by the Perron eigenvector of the underlying primitive morphism, namely sqrt(3)/2 + 3/2 $\approx$ 2.366025. - \char95{}Joaquim Manji-Jung\char95{}, Nov 19 2026\par
\%C A284492 The conjectured limit is incorrect: a(n)/n tends to 1/f, where f is the frequency of the letter given by the Perron eigenvector of the underlying primitive morphism, namely sqrt(3) $\approx$ 1.732051. - \char95{}Joaquim Manji-Jung\char95{}, Nov 19 2026\par
\%C A285413 The conjectured limit is incorrect: a(n)/n tends to 1/f, where f is the frequency of the letter given by the Perron eigenvector of the underlying primitive morphism, namely 3/4 + sqrt(17)/4 $\approx$ 1.780776. - \char95{}Joaquim Manji-Jung\char95{}, Nov 19 2026\par
\%C A285428 The conjectured limit is incorrect: a(n)/n tends to 1/f, where f is the frequency of the letter given by the Perron eigenvector of the underlying primitive morphism, namely sqrt(17)/4 + 5/4 $\approx$ 2.280776. - \char95{}Joaquim Manji-Jung\char95{}, Nov 19 2026\par
\%C A285429 The conjectured limit is incorrect: a(n)/n tends to 1/f, where f is the frequency of the letter given by the Perron eigenvector of the underlying primitive morphism, namely 3/4 + sqrt(17)/4 $\approx$ 1.780776. - \char95{}Joaquim Manji-Jung\char95{}, Nov 19 2026\par
\%C A284684 The conjecture is true: the letter frequency of this transform of the Thue-Morse word is given exactly by the frequencies of factors of A010060, which are rational and computable from the Perron eigenvector of the induced block substitution (see the companion repository). - \char95{}Joaquim Manji-Jung\char95{}, Nov 19 2026\par
\%C A284685 The conjecture is true: the letter frequency of this transform of the Thue-Morse word is given exactly by the frequencies of factors of A010060, which are rational and computable from the Perron eigenvector of the induced block substitution (see the companion repository). - \char95{}Joaquim Manji-Jung\char95{}, Nov 19 2026\par
\%C A284790 The conjecture is true: the letter frequency of this transform of the Thue-Morse word is given exactly by the frequencies of factors of A010060, which are rational and computable from the Perron eigenvector of the induced block substitution (see the companion repository). - \char95{}Joaquim Manji-Jung\char95{}, Nov 19 2026\par
\%C A284791 The conjecture is true: the letter frequency of this transform of the Thue-Morse word is given exactly by the frequencies of factors of A010060, which are rational and computable from the Perron eigenvector of the induced block substitution (see the companion repository). - \char95{}Joaquim Manji-Jung\char95{}, Nov 19 2026\par
\%C A284933 The conjecture is true: the letter frequency of this transform of the Thue-Morse word is given exactly by the frequencies of factors of A010060, which are rational and computable from the Perron eigenvector of the induced block substitution (see the companion repository). - \char95{}Joaquim Manji-Jung\char95{}, Nov 19 2026\par
\%C A284934 The conjecture is true: the letter frequency of this transform of the Thue-Morse word is given exactly by the frequencies of factors of A010060, which are rational and computable from the Perron eigenvector of the induced block substitution (see the companion repository). - \char95{}Joaquim Manji-Jung\char95{}, Nov 19 2026\par
\%C A285091 The conjecture is true: the letter frequency of this transform of the Thue-Morse word is given exactly by the frequencies of factors of A010060, which are rational and computable from the Perron eigenvector of the induced block substitution (see the companion repository). - \char95{}Joaquim Manji-Jung\char95{}, Nov 19 2026\par
\%C A285092 The conjecture is true: the letter frequency of this transform of the Thue-Morse word is given exactly by the frequencies of factors of A010060, which are rational and computable from the Perron eigenvector of the induced block substitution (see the companion repository). - \char95{}Joaquim Manji-Jung\char95{}, Nov 19 2026\par
\%C A285958 The conjecture is true: the letter frequency of this transform of the Thue-Morse word is given exactly by the frequencies of factors of A010060, which are rational and computable from the Perron eigenvector of the induced block substitution (see the companion repository). - \char95{}Joaquim Manji-Jung\char95{}, Nov 19 2026\par
\%C A285959 The conjecture is true: the letter frequency of this transform of the Thue-Morse word is given exactly by the frequencies of factors of A010060, which are rational and computable from the Perron eigenvector of the induced block substitution (see the companion repository). - \char95{}Joaquim Manji-Jung\char95{}, Nov 19 2026\par
\%C A285965 The conjecture is true: the letter frequency of this transform of the Thue-Morse word is given exactly by the frequencies of factors of A010060, which are rational and computable from the Perron eigenvector of the induced block substitution (see the companion repository). - \char95{}Joaquim Manji-Jung\char95{}, Nov 19 2026\par
\%C A285973 The conjecture is true: the letter frequency of this transform of the Thue-Morse word is given exactly by the frequencies of factors of A010060, which are rational and computable from the Perron eigenvector of the induced block substitution (see the companion repository). - \char95{}Joaquim Manji-Jung\char95{}, Nov 19 2026\par
\%C A285974 The conjecture is true: the letter frequency of this transform of the Thue-Morse word is given exactly by the frequencies of factors of A010060, which are rational and computable from the Perron eigenvector of the induced block substitution (see the companion repository). - \char95{}Joaquim Manji-Jung\char95{}, Nov 19 2026\par
\%C A285976 The conjecture is true: the letter frequency of this transform of the Thue-Morse word is given exactly by the frequencies of factors of A010060, which are rational and computable from the Perron eigenvector of the induced block substitution (see the companion repository). - \char95{}Joaquim Manji-Jung\char95{}, Nov 19 2026\par
\%C A285979 The conjecture is true: the letter frequency of this transform of the Thue-Morse word is given exactly by the frequencies of factors of A010060, which are rational and computable from the Perron eigenvector of the induced block substitution (see the companion repository). - \char95{}Joaquim Manji-Jung\char95{}, Nov 19 2026\par
\%C A285980 The conjecture is true: the letter frequency of this transform of the Thue-Morse word is given exactly by the frequencies of factors of A010060, which are rational and computable from the Perron eigenvector of the induced block substitution (see the companion repository). - \char95{}Joaquim Manji-Jung\char95{}, Nov 19 2026\par
\%C A286045 The conjecture is true: the letter frequency of this transform of the Thue-Morse word is given exactly by the frequencies of factors of A010060, which are rational and computable from the Perron eigenvector of the induced block substitution (see the companion repository). - \char95{}Joaquim Manji-Jung\char95{}, Nov 19 2026\par
\%C A286053 The conjecture is true: the letter frequency of this transform of the Thue-Morse word is given exactly by the frequencies of factors of A010060, which are rational and computable from the Perron eigenvector of the induced block substitution (see the companion repository). - \char95{}Joaquim Manji-Jung\char95{}, Nov 19 2026\par
\%C A003157 The conjecture is true: the letter frequency of this transform of the Thue-Morse word is given exactly by the frequencies of factors of A010060, which are rational and computable from the Perron eigenvector of the induced block substitution (see the companion repository). - \char95{}Joaquim Manji-Jung\char95{}, Nov 19 2026\par
\%C A285964 The limit is 5/4, not 4/5 (which is impossible since a(n) \char62{}= n): the frequency of 0 in A285963 is 4/5. - \char95{}Joaquim Manji-Jung\char95{}, Nov 19 2026\par
\%C A285590 The conjecture is true: the transform halves (rounding up) each maximal run of the underlying morphic word, whose run frequencies are computed exactly from the Perron eigenvector of the induced block substitution (see the companion repository). - \char95{}Joaquim Manji-Jung\char95{}, Nov 19 2026\par
\%C A285591 The conjecture is true: the transform halves (rounding up) each maximal run of the underlying morphic word, whose run frequencies are computed exactly from the Perron eigenvector of the induced block substitution (see the companion repository). - \char95{}Joaquim Manji-Jung\char95{}, Nov 19 2026\par
\%C A285593 The conjecture is true: the transform halves (rounding up) each maximal run of the underlying morphic word, whose run frequencies are computed exactly from the Perron eigenvector of the induced block substitution (see the companion repository). - \char95{}Joaquim Manji-Jung\char95{}, Nov 19 2026\par
\%C A285594 The conjecture is true: the transform halves (rounding up) each maximal run of the underlying morphic word, whose run frequencies are computed exactly from the Perron eigenvector of the induced block substitution (see the companion repository). - \char95{}Joaquim Manji-Jung\char95{}, Nov 19 2026\par
}
\subsection*{\texttt{conjectures/precalg/draft.txt}}
{\small\raggedright\ttfamily
\%C A025258 The conjectured recurrence is true (for n \char62{}= 6): with theta = x d/dx it amounts to L(A) being a polynomial for L = Sum\char95{}i x\char94{}i p\char95{}i(theta+i), and reducing L(A) modulo the algebraic equation -A\char94{}2 - A*x\char94{}2 + A*x + A - x\char94{}3 + x\char94{}2 - x = 0 of the g.f. (Y = A(x), Y' = -P\char95{}x/P\char95{}Y) indeed gives a polynomial. It fails at n = 5. - \char95{}Joaquim Manji-Jung\char95{}, Nov 19 2026\par
\%C A059279 The conjectured recurrence is true (for n \char62{}= 3): with theta = x d/dx it amounts to L(A) being a polynomial for L = Sum\char95{}i x\char94{}i p\char95{}i(theta+i), and reducing L(A) modulo the algebraic equation 2*A\char94{}2*x\char94{}2 - A\char94{}2*x - 2*A*x + A + x - 1 = 0 of the g.f. (Y = A(x), Y' = -P\char95{}x/P\char95{}Y) indeed gives a polynomial. - \char95{}Joaquim Manji-Jung\char95{}, Nov 19 2026\par
\%C A073155 The conjectured recurrence is true (for n \char62{}= 2): with theta = x d/dx it amounts to L(A) being a polynomial for L = Sum\char95{}i x\char94{}i p\char95{}i(theta+i), and reducing L(A) modulo the algebraic equation -A\char94{}2*x\char94{}3 - 2*A\char94{}2*x\char94{}2 - A\char94{}2*x + A - 1 = 0 of the g.f. (Y = A(x), Y' = -P\char95{}x/P\char95{}Y) indeed gives a polynomial. - \char95{}Joaquim Manji-Jung\char95{}, Nov 19 2026\par
\%C A080243 The conjectured recurrence is true (for n \char62{}= 2): with theta = x d/dx it amounts to L(A) being a polynomial for L = Sum\char95{}i x\char94{}i p\char95{}i(theta+i), and reducing L(A) modulo the algebraic equation 2*A\char94{}2*x - A*x + A - 1 = 0 of the g.f. (Y = A(x), Y' = -P\char95{}x/P\char95{}Y) indeed gives a polynomial. - \char95{}Joaquim Manji-Jung\char95{}, Nov 19 2026\par
\%C A114584 The conjectured recurrence is true (for n \char62{}= 5): with theta = x d/dx it amounts to L(A) being a polynomial for L = Sum\char95{}i x\char94{}i p\char95{}i(theta+i), and reducing L(A) modulo the algebraic equation -A\char94{}2*x\char94{}2 + A*x\char94{}3 - A*x + A - 1 = 0 of the g.f. (Y = A(x), Y' = -P\char95{}x/P\char95{}Y) indeed gives a polynomial. - \char95{}Joaquim Manji-Jung\char95{}, Nov 19 2026\par
\%C A143013 The conjectured recurrence is true (for n \char62{}= 3): with theta = x d/dx it amounts to L(A) being a polynomial for L = Sum\char95{}i x\char94{}i p\char95{}i(theta+i), and reducing L(A) modulo the algebraic equation -A\char94{}2*x\char94{}2 - A*x + A - x - 1 = 0 of the g.f. (Y = A(x), Y' = -P\char95{}x/P\char95{}Y) indeed gives a polynomial. - \char95{}Joaquim Manji-Jung\char95{}, Nov 19 2026\par
\%C A166300 The conjectured recurrence is true (for n \char62{}= 3): with theta = x d/dx it amounts to L(A) being a polynomial for L = Sum\char95{}i x\char94{}i p\char95{}i(theta+i), and reducing L(A) modulo the algebraic equation A\char94{}2*x\char94{}4 - A\char94{}2*x - A*x\char94{}3 + A + x - 1 = 0 of the g.f. (Y = A(x), Y' = -P\char95{}x/P\char95{}Y) indeed gives a polynomial. - \char95{}Joaquim Manji-Jung\char95{}, Nov 19 2026\par
\%C A174169 The conjectured recurrence is true (for n \char62{}= 3): with theta = x d/dx it amounts to L(A) being a polynomial for L = Sum\char95{}i x\char94{}i p\char95{}i(theta+i), and reducing L(A) modulo the algebraic equation -A\char94{}2*x\char94{}2 + 3*A*x\char94{}2 - A*x + A - 1 = 0 of the g.f. (Y = A(x), Y' = -P\char95{}x/P\char95{}Y) indeed gives a polynomial. - \char95{}Joaquim Manji-Jung\char95{}, Nov 19 2026\par
\%C A174171 The conjectured recurrence is true (for n \char62{}= 3): with theta = x d/dx it amounts to L(A) being a polynomial for L = Sum\char95{}i x\char94{}i p\char95{}i(theta+i), and reducing L(A) modulo the algebraic equation -A\char94{}2*x\char94{}2 - 2*A*x\char94{}2 - A*x + A - 1 = 0 of the g.f. (Y = A(x), Y' = -P\char95{}x/P\char95{}Y) indeed gives a polynomial. - \char95{}Joaquim Manji-Jung\char95{}, Nov 19 2026\par
\%C A188314 The conjectured recurrence is true (for n \char62{}= 5): with theta = x d/dx it amounts to L(A) being a polynomial for L = Sum\char95{}i x\char94{}i p\char95{}i(theta+i), and reducing L(A) modulo the algebraic equation -A\char94{}2*x - A*x\char94{}2 + A - x - 1 = 0 of the g.f. (Y = A(x), Y' = -P\char95{}x/P\char95{}Y) indeed gives a polynomial. - \char95{}Joaquim Manji-Jung\char95{}, Nov 19 2026\par
\%C A215973 The conjectured recurrence is true (for n \char62{}= 4): with theta = x d/dx it amounts to L(A) being a polynomial for L = Sum\char95{}i x\char94{}i p\char95{}i(theta+i), and reducing L(A) modulo the algebraic equation A\char94{}2*x\char94{}2 - A\char94{}2*x - 2*A*x + A + x - 1 = 0 of the g.f. (Y = A(x), Y' = -P\char95{}x/P\char95{}Y) indeed gives a polynomial. - \char95{}Joaquim Manji-Jung\char95{}, Nov 19 2026\par
\%C A217333 The conjectured recurrence is true (for n \char62{}= 3): with theta = x d/dx it amounts to L(A) being a polynomial for L = Sum\char95{}i x\char94{}i p\char95{}i(theta+i), and reducing L(A) modulo the algebraic equation -A\char94{}2*x\char94{}3 - 2*A*x + A + x - 1 = 0 of the g.f. (Y = A(x), Y' = -P\char95{}x/P\char95{}Y) indeed gives a polynomial. - \char95{}Joaquim Manji-Jung\char95{}, Nov 19 2026\par
\%C A254314 The conjectured recurrence is true (for n \char62{}= 5): with theta = x d/dx it amounts to L(A) being a polynomial for L = Sum\char95{}i x\char94{}i p\char95{}i(theta+i), and reducing L(A) modulo the algebraic equation -2*A\char94{}2*x + A\char94{}2 - A*x\char94{}2 + 6*A*x - 3*A - 3*x + 2 = 0 of the g.f. (Y = A(x), Y' = -P\char95{}x/P\char95{}Y) indeed gives a polynomial. - \char95{}Joaquim Manji-Jung\char95{}, Nov 19 2026\par
\%C A257178 The conjectured recurrence is true (for n \char62{}= 2): with theta = x d/dx it amounts to L(A) being a polynomial for L = Sum\char95{}i x\char94{}i p\char95{}i(theta+i), and reducing L(A) modulo the algebraic equation 3*A\char94{}2*x\char94{}3 - A\char94{}2*x\char94{}2 - 3*A*x + A - 1 = 0 of the g.f. (Y = A(x), Y' = -P\char95{}x/P\char95{}Y) indeed gives a polynomial. - \char95{}Joaquim Manji-Jung\char95{}, Nov 19 2026\par
\%C A257388 The conjectured recurrence is true (for n \char62{}= 2): with theta = x d/dx it amounts to L(A) being a polynomial for L = Sum\char95{}i x\char94{}i p\char95{}i(theta+i), and reducing L(A) modulo the algebraic equation 4*A\char94{}2*x\char94{}3 - A\char94{}2*x\char94{}2 - 4*A*x + A - 1 = 0 of the g.f. (Y = A(x), Y' = -P\char95{}x/P\char95{}Y) indeed gives a polynomial. - \char95{}Joaquim Manji-Jung\char95{}, Nov 19 2026\par
\%C A261196 The conjectured recurrence is true (for n \char62{}= 0): with theta = x d/dx it amounts to L(A) being a polynomial for L = Sum\char95{}i x\char94{}i p\char95{}i(theta+i), and reducing L(A) modulo the algebraic equation A\char94{}4 - 16*A\char94{}2*x - 1 = 0 of the g.f. (Y = A(x), Y' = -P\char95{}x/P\char95{}Y) indeed gives a polynomial. - \char95{}Joaquim Manji-Jung\char95{}, Nov 19 2026\par
\%C A346074 The conjectured recurrence is true (for n \char62{}= 2): with theta = x d/dx it amounts to L(A) being a polynomial for L = Sum\char95{}i x\char94{}i p\char95{}i(theta+i), and reducing L(A) modulo the algebraic equation A\char94{}2*x\char94{}6 - A\char94{}2*x\char94{}5 - A*x + A - 1 = 0 of the g.f. (Y = A(x), Y' = -P\char95{}x/P\char95{}Y) indeed gives a polynomial. - \char95{}Joaquim Manji-Jung\char95{}, Nov 19 2026\par
}
\subsection*{\texttt{conjectures/yanev/draft.txt}}
{\small\raggedright\ttfamily
\%C A129527 Yanev's conjecture a(n) = 2n - A000265(n) is true: the right side satisfies f(2n+1) = 2n+1 and f(2n) = f(n) + 2n. - \char95{}Joaquim Manji-Jung\char95{}, Nov 19 2026\par
\%C A173614 Yanev's conjecture a(n) = A002322(A007947(n)) is true: for squarefree m, lambda(m) = lcm\char95{}\char123{}p|m\char125{} lambda(p) = lcm\char95{}\char123{}p|m\char125{} (p-1). - \char95{}Joaquim Manji-Jung\char95{}, Nov 19 2026\par
\%C A023887 Yanev's conjecture is true: both sides are multiplicative and agree at prime powers, sigma(p\char94{}(me+m-1))/sigma(p\char94{}(m-1)) = (p\char94{}(m(e+1))-1)/(p\char94{}m-1). - \char95{}Joaquim Manji-Jung\char95{}, Nov 19 2026\par
\%C A055210 Yanev's conjecture is true: both sides are multiplicative and equal (p\char94{}(2t+1)+1)/(p+1) at n = p\char94{}e, t = floor(e/2). - \char95{}Joaquim Manji-Jung\char95{}, Nov 19 2026\par
\%C A367865 Yanev's conjecture is true: both sides equal Product\char95{}\char123{}p|n\char125{} (p\char94{}2-p+1), the right side because rad(nk)/rad(k) = Product\char95{}\char123{}p|n, p does not divide k\char125{} p depends only on k mod rad(n). - \char95{}Joaquim Manji-Jung\char95{}, Nov 19 2026\par
\%C A090079 Yanev's conjecture is true: a(n) is the alternating word 1010... of length r = A005811(n), whose value is (2\char94{}(r+1)-1)/3 if it ends in 1 (n odd) and (2\char94{}(r+1)-2)/3 otherwise. - \char95{}Joaquim Manji-Jung\char95{}, Nov 19 2026\par
\%C A163617 Yanev's conjecture is true: using A003188(2n) = 2*A003188(n) + (n mod 2) and A003188(2n+1) = 2*A003188(n) + 1 - (n mod 2), the right side satisfies the defining recurrences. - \char95{}Joaquim Manji-Jung\char95{}, Nov 19 2026\par
\%C A285385 Yanev's conjecture is true: the terms are the m such that m-1 has an even number of trailing zeros in binary (see A072939), so the next term after m is m+1 if A007814(m) is even and m+2 otherwise, and A096268(m-1) = A007814(m) mod 2. - \char95{}Joaquim Manji-Jung\char95{}, Nov 19 2026\par
\%C A179504 Yanev's conjecture is true: a(n)/A023887(n) = sigma\char95{}n(2\char94{}(v+1))/sigma\char95{}n(2\char94{}v) = (2\char94{}(n(v+2))-1)/(2\char94{}(n(v+1))-1), v = A007814(n), which equals the stated factor. - \char95{}Joaquim Manji-Jung\char95{}, Nov 19 2026\par
}
